\documentclass[11pt]{article}
\usepackage[a4paper,margin=1in]{geometry}
\usepackage[utf8]{inputenc}
\usepackage[T1]{fontenc}
\usepackage{parskip}
\usepackage{amsmath,amssymb,amsthm,mathtools}
\usepackage{booktabs,longtable,array}
\usepackage{etoolbox}
% Local compatibility with the installed pre-2026-Q2 longtable output routine.
% Finite-shrink correction: https://github.com/latex3/latex2e/issues/1907
\makeatletter
\patchcmd{\LT@output}{\vss}{\vskip 0pt plus \maxdimen minus \normalbaselineskip}{}{}
\patchcmd{\LT@output}{\vss}{\vskip 0pt plus \maxdimen minus \normalbaselineskip}{}{}
\makeatother
\usepackage{graphicx,xcolor}
\usepackage{enumitem}
\usepackage{fancyvrb}
\usepackage[hidelinks]{hyperref}
\setlength{\parindent}{0pt}
\setlength{\emergencystretch}{2em}
\setcounter{LTchunksize}{100}
\setlist{itemsep=3pt,topsep=4pt}
\allowdisplaybreaks[1]
\newcolumntype{L}[1]{>{\raggedright\arraybackslash}p{#1}}
\newtheorem{proposition}{Proposition}[section]
\newtheorem{theorem}[proposition]{Theorem}
\newtheorem{definition}[proposition]{Definition}
\theoremstyle{remark}
\newtheorem{remark}[proposition]{Scope}
\newcommand{\Tr}{\operatorname{Tr}}
\newcommand{\Inv}{\operatorname{Inv}}
\newcommand{\Ran}{\operatorname{Ran}}
\newcommand{\Cov}{\operatorname{Cov}}
\newcommand{\supp}{\operatorname{supp}}
\newcommand{\R}{\mathbb R}
\newcommand{\C}{\mathbb C}
\newcommand{\Gstar}{G_*}
\newcommand{\Lstar}{L_*}
\newcommand{\tstar}{\tau_*}
\newcommand{\Ach}{A_{\rm ch}}
\newcommand{\eps}{\epsilon_e}
\numberwithin{equation}{section}
\hypersetup{pdftitle={A microscopic realization of causal triangulations with native matter and clock dynamics},pdfauthor={Jacob Chinitz},pdfsubject={Exact constructions, Einstein normalization, conditional phase inheritance, and replication}}
\author{Jacob Chinitz}
\title{A microscopic realization of causal triangulations with native matter and clock dynamics\\[0.3em]
\large Capacity construction, CDT weights, Einstein normalization, and reproducible finite tests}
\date{October 20, 2026}

\begin{document}
\maketitle

\begin{abstract}
A finite entanglement-capacity construction is extended to an explicit microscopic realization of a selected causal triangulation measure. The local $1680$-state ensemble retains its closure spectrum, entropy, and charged rest gap. An equivariant encoding exposes four coherent binary bond modes inside the existing capacity states. A specified absolute angular transfer converts their cut entropy into a Regge curvature weight. Connected spin-three geometric patches, normalized by their full Gram matrices, realize all Euclidean tetrahedral interface spaces of an $\alpha=7$ causal regulator. Their tensor contraction equals the selected CDT history weight, including boundary terms and compatible external sources. A paired-cell event-volume calibration gives $G_0/\Gstar=0.10551926870131369\,m_4^{8/3}$, or $1.03002098319$ at the provisional chain constant $m_4=2.35$.

A second construction couples an enumerated two-geometry transfer to retained orbital cochains. Harmonic-preserving polar transport keeps the native rest gap exact after geometry coupling. The mass derivative of the joint Hamiltonian defines a positive integrated proper clock. A cochain retraction yields a positive microscopic metric that intertwines the entire supplied coarse generator, its clock, and transported source operators; the joint intertwining residuals are a few $10^{-13}$ or smaller, while the separate finite-difference clock checks are at the $10^{-10}$ level. The manuscript's Einstein matching, this bare normalization, and an interacting response measurement are compared explicitly. An inheritance theorem states the conditions under which de Sitter observables and a renormalized gravitational response transfer from CDT to the microscopic realization. The finite construction satisfies the algebraic correspondence. At its fixed $\alpha=7$ curvature couplings, the registered volume-stiffness ratio is $G_{\rm vol}/G_0=1.064$, with logarithmic standard error $0.089$, inside the $0.80$--$1.25$ agreement band. This primary result uses one stationary 32k run, with reported-only stationary runs at 16k and 64k supporting the comparison. Start independence remains undetermined. A separate curvature-trace test is inconclusive, and exploratory scalar binding has no stable energy extraction. A common interacting matter cone and the long-distance source response remain open.
\end{abstract}

\clearpage
{\small\tableofcontents}
\clearpage
\paragraph{Scope and companion papers.}
Prepared for Jacob Chinitz as a continuation of \emph{Entropic Scalar EFT: From Entanglement Microstructure to Gravity and Cosmic Structure}.

The construction starts with the finite capacity ensemble and its retained charged sector. It specifies the additional geometric interactions, proves an exact correspondence with a selected Euclidean CDT state sum, and constructs a finite joint geometry--matter--clock operator. The accompanying replication bundle contains the inputs, executable calculations, reference outputs, and verification records.

The de Sitter and interacting Einstein conclusions have explicit hypotheses: the same coarse ensemble must occupy the required phase, and its physical source response must have the appropriate infrared limit. The exact microscopic correspondence transfers such results when those hypotheses hold.
The current Astra manuscript incorporates this geometric and matter--clock construction. The connected ensemble has registered finite-volume stiffness agreement, with its matter-source Newton coefficient still unresolved. The comparison and evidence sections state these distinct physical results.

\section{Physical variables and the meaning of a microscopic realization}
\label{sec:premises}

The microscopic substrate carries a finite capacity for entanglement, localized retained defects, and relational geometric data. The capacity variables describe which finite states are available and how a defect is maintained through renewal. The geometric variables describe areas, shapes, and parallel transport. Their roles are fixed before the continuum approximation is taken.

The ordinary gravitational interpretation uses one metric. In a static branch, a capacity deficit can be a function of its lapse and boundary normalization. A second freely propagating scalar is unnecessary. The Einstein action, expressed in that capacity variable, supplies the ordinary Newton and metric response. A microscopic construction must therefore reproduce both the coefficient and the operator structure of the metric response.

The local capacity count alone leaves geometric dynamics undetermined. A Hamiltonian can have the same finite local state space and many different interactions. The selected realization fixes those interactions by explicit matrices and tensor contractions. Its mathematical content can be evaluated without requiring a uniqueness theorem for all admissible microscopic actions.

\subsection{Three gravitational coefficients}
\label{sec:threeG}

Use $\hbar=c=1$ unless SI factors are displayed. A proper tick has duration $\tstar$ and causal length $\Lstar=c\tstar$. The three gravitational coefficients are
\begin{align}
\Gstar&=\frac{c^3\Lstar^2}{\hbar},\\
G_0&=\text{coefficient in the selected microscopic curvature weight},\\
G_R&=\text{coefficient extracted from an interacting long-distance response}.
\end{align}
The electron-calibrated $\Gstar$ is a dimensional reference. The ordinary matched branch of the main manuscript selects $G=\Gstar$; the empirical companion uses the usual measured Newton constant in its effective astrophysical equations. Neither use is a measurement of this new ensemble's $G_R$. The curvature construction below determines $G_0/\Gstar$ under its stated volume and angular prescriptions. A measurement of $G_R$ additionally involves the full geometry measure, the source, and a physical infrared regime.

An exact microscopic representation of a coarse partition function preserves all observables included in that representation. It also preserves the coarse theory's dependence on couplings and phase. Thus a microscopic representation of a CDT model in a de Sitter regime inherits its de Sitter observables. A representation at another coupling point inherits the observables at that point. The correspondence does not identify distinct coupling points.

\subsection{The selected data}

The capacity part inherits the following commitments from the main manuscript: the tetrahedral architecture, positive oriented face matching, injective occupation of four relational ports, orientation doubling, faithful renewal of the replaceable capacity factor, and the selected marked charged generator. The electron provides the dimensional anchor.

The continuation supplies a coherent global bond state, an absolute angular-transfer rule, a paired-cell volume calibration, and connected geometric interface maps. It selects causal triangulations as the coarse history space. Its matter extension selects a positive cochain metric, a positive-energy orbital generator, and normalized transport between different spatial geometries.

Each supplied item has a distinct test. The local Gibbs state can be enumerated exactly. The geometric maps can be checked for rank, isometry, and gauge covariance. The state-sum equality is a tensor-contraction identity. The physical phase and the long-distance response are properties of the summed interacting measure.

\subsection{Notation and operator domains}

\begin{longtable}{L{0.20\textwidth}L{0.74\textwidth}}
\toprule Symbol & Meaning\\\midrule\endhead
$\mathcal H_{\rm cap,o}$ & The $1680$-dimensional capacity and orientation space.\\
$\mathcal I_j=\Inv(V_j^{\otimes4})$ & Gauge-invariant geometric tetrahedron space with equal face spins $j$.\\
$Q_{\rm mis}$ & The retained nine-dimensional mismatch fiber; its projector is written $P_{\rm mis}$.\\
$K^2$ & The finite capacity closure statistic. It is distinct from a geometric Gauss Casimir.\\
$g_{\rm share}$ & Entropy of the complete local capacity and orientation state.\\
$\eps$ & Dimensionless electron rest gap per proper tick.\\
$m_4$ & Four-dimensional asymptotic longest-chain constant in the specified event law.\\
$v_*$ & Physical four-volume per clock event. In $c=1$ units it has length dimension four.\\
$a,\alpha$ & CDT spatial edge and Euclidean mixed-edge squared ratio.\\
$N_h$ & Number of effective binary channels assigned to hinge $h$.\\
$N$ & Lapse or, in the half-transfer theorem, a supplied positive duration operator; the role is stated locally.\\
$P,R,J$ & Cochain interpolation, cochain retraction, and its orthonormal isometric encoding.\\
$\Delta_{\rm CDT},\Delta t$ & Bare CDT asymmetry coupling and reference time interval.\\
\bottomrule
\end{longtable}

Across the three accounts, $g_{\rm share}$ here is Astra's
$g_{\rm share,eff}$ and the empirical paper's $g$.
The electron gap $\eps$ is distinct from the galactic loading
$\epsilon=g/(4\pi^2)$. The physical event volume $v_*$ here is
$v_{\rm ev}$ in Astra C.8; neither is automatically $\Lstar^4$,
the volume assigned in the earlier canonical reference convention.

Finite-dimensional calculations use the positive Hilbert-space inner products stated with each operator. A Euclidean transfer is a positive operator when its quadratic form is nonnegative. This property is stronger than positivity of selected observables and different from entrywise nonnegativity of a matrix in one basis. The logarithm of a transfer is taken on its nonzero support.

\section{The native finite capacity ensemble}
\label{sec:capacity}

\subsection{The face split and the seven-letter branch}

Two primitive spin-$3/2$ slots have
\begin{equation}
V_{3/2}\otimes V_{3/2}=V_0\oplus V_1\oplus V_2\oplus V_3.
\end{equation}
The selected positive mismatch operator is
\begin{equation}
Q_f=6P_0+5P_1+3P_2,\qquad \ker Q_f=V_3.
\label{eq:faceQ}
\end{equation}
The transmitted sector has dimension seven; its orthogonal complement has dimension $1+3+5=9$. The sharp reversible realization identifies that complete complement with the marked mismatch carrier. This identification is part of the selected microscopic vertex. Given it, the $7+9$ split is exact.

The native capacity alphabet is $m=-3,-2,\ldots,3$. A cell chooses four distinct letters, assigns them to four ordered ports, and carries one of two orientations. Therefore
\begin{equation}
\dim\mathcal H_{\rm cap,o}
=\binom74\,4!\,2=1680.
\label{eq:count}
\end{equation}
Injectivity is imposed within a cell. The graph-level rules must specify separately how neighboring cells share or transport their data.

In the selected relational geometry, a primitive cell also carries
\begin{equation}
\mathcal I_3=\Inv_{SU(2)}(V_3^{\otimes4}),\qquad \dim\mathcal I_3=7.
\end{equation}
The dimension follows by coupling the first pair and second pair to the same intermediate spin $k=0,\ldots,6$, then coupling those two spins to zero. This geometric factor remains distinct from the capacity alphabet. Their equal dimensions do not identify their operators.

\subsection{Closure weights and entropy}

For an unordered four-letter set $s$, define
\begin{equation}
K^2(s)=48-\frac{(\sum_{m\in s}m)^2-\sum_{m\in s}m^2}{3}.
\label{eq:K2}
\end{equation}
Routing and orientation leave $K^2$ unchanged. The admissibility law is
\begin{equation}
P_s(\eta)=\frac{e^{-\eta K^2(s)}}{\sum_{s'}e^{-\eta K^2(s')}}.
\end{equation}
Every routing-orientation state above $s$ has probability $P_s/48$. The selected three-dimensional closure condition fixes $\eta$ through
\begin{equation}
\eta\langle K^2\rangle_\eta=\frac32.
\label{eq:closure}
\end{equation}
There are only $35$ unordered sets. Enumeration followed by a scalar root solve gives
\begin{align}
\eta_*&=0.0298668443935226,\\
g_{\rm share}
&=-\sum_sP_s\log P_s+\log48
=7.4198000235709.
\label{eq:entropy}
\end{align}
The same computation can enumerate the $840$ ordered distinct tuples and add $\log2$ for orientation. Both procedures must give the same entropy. They provide an immediate bookkeeping test of the implementation.

The closure condition, finite alphabet, and state measure define the numerical problem. A different measure or relaxation of injectivity defines a different ensemble. The constants above are consequences of the specified ensemble.

\subsection{Faithful renewal}

For a stationary process with one-step marginal $p_*$,
\begin{equation}
H(B_{t+1}\mid B_t)=H(p_*)-I(B_t;B_{t+1}).
\end{equation}
Faithful renewal requires each pass to carry the full marginal entropy. It therefore sets the endpoint mutual information to zero. The replacement kernel is
\begin{equation}
\mathsf K(b,b')=p_*(b').
\label{eq:refresh}
\end{equation}
This proves independence of successive replaceable capacity marks within the chosen class. A law for spatial event positions remains a separate part of the clock construction.

A reversible implementation exports the old register into history and prepares a fresh register with the prescribed state. The global dilation is isometric. The reduced present register follows a replacement channel after the exported degrees of freedom are traced out. This division of retained and replaced factors allows a particle to persist while its surrounding capacity renews.

\subsection{The charged rest generator}

The marked generator uses the recurrence weight $r=e^{-7g_{\rm share}}$ and the selected finite correction
\begin{align}
z&=9e^{-g_{\rm share}}
\left(1-\frac{8\eta_*}{49}\right)^{21/2},\\
Z_e&=(1+z)(1+7z^2),\\
\eps&=-\frac32 Z_e\log(1-r)
=4.185458795583485\times10^{-23}.
\label{eq:epsilon}
\end{align}
These formulas include the marked routing and additive-generator prescription of Appendix H.9. The local state count supplies the entropy and recurrence weight; the prescription specifies how the weights enter the physical generator.

For an electron occupation projector $P_e$ the rest transfer is
\begin{equation}
T_e(s)=e^{-\eps sP_e}.
\end{equation}
The electron calibration reads
\begin{equation}
m_ec^2=\frac{\hbar\eps}{\tstar},\qquad
\Lstar=c\tstar=\frac{\hbar\eps}{m_ec}.
\label{eq:electronanchor}
\end{equation}
It fixes the proper cadence and its causal length. A spatial hopping coefficient must be checked against this same $c$. The number $\eps$ should be evaluated with a stable function such as \texttt{log1p}; direct subtraction of $1-e^{-7g_{\rm share}}$ at machine precision can erase its small difference from one.

The local exclusion cost $\log(7/6)$, the elementary binary cut entropy $\log2$, the maintenance throughput, and the rest gap have different definitions. A persistent particle's inertial energy is read from the retained transfer. Substituting a local likelihood cost for that energy changes the source.

\subsection{The manuscript's static normalization}
\label{sec:mainnormalization}

The manuscript's diamond-graph scalar calculation uses the unnormalized Laplacian $L_\diamond=4I-A$. Its on-site inverse is
\begin{equation}
G_{\rm tet}(0)
=\frac{3\Gamma(1/3)^6}{2^{14/3}\pi^4}
=0.4482203943883814\ldots .
\end{equation}
With each undirected bond counted once, the stiffness and source amplitudes share the same edge coupling. Their ratio is
\begin{equation}
\frac{\kappa}{\gamma}
=\frac{3\Lstar}{4G_{\rm tet}(0)\kappa_m},
\qquad
\kappa_m=\frac{\hbar}{c\Lstar\log2}.
\end{equation}
The weak-field bridge gives
\begin{equation}
\frac{G}{\Gstar}=\frac{S_0}{S_\infty},\qquad
S_0=\frac{3\log2}{32\pi G_{\rm tet}(0)}
=0.0461482516\ldots .
\label{eq:S0}
\end{equation}
Appendix C.5 selects the physical matching $S_\infty=S_0$, hence $G=\Gstar$. An entropy-unit change leaves $\kappa/(\gamma S_\infty)$ invariant. A change in the physical matching at fixed entropy convention changes the inferred $G$.

The continuation below preserves the native finite constants and supplies a different explicit route to a bare geometric coefficient through event volume and angular transfer. Comparing the two routes requires carrying their physical normalization premises along with the arithmetic.

\section{Coherent bonds in the existing capacity states}
\label{sec:bonds}

\subsection{Equivariant factorization of routing and orientation}

For each unordered four-letter set, routing carries the regular representation $\C[S_4]$. Orientation supplies an additional two-dimensional multiplicity space. The routing-orientation representation admits
\begin{equation}
\C[S_4]\otimes\C^2
\cong\mathcal F(\C^4)\otimes V_{\rm std},
\label{eq:fockidentity}
\end{equation}
where $\mathcal F(\C^4)$ is the fermionic Fock space of four virtual modes and $V_{\rm std}$ is the three-dimensional standard representation of $S_4$.

The identity follows from characters. In the order identity, transposition, double transposition, three-cycle, four-cycle, the Fock character is $(16,0,0,4,0)$ and the standard character is $(3,1,-1,0,-1)$. Their product is $(48,0,0,0,0)$, the character of two regular representations. Equality of finite-group characters proves unitary equivalence.

Including the $35$ unordered sets gives
\begin{equation}
\mathcal H_{\rm cap,o}\cong\mathcal F_4\otimes\C^{105}.
\end{equation}
The virtual occupations are observables encoded in the existing routing and orientation states. The local Hilbert dimension remains $1680$. The fermionic signs in the port-permutation action are required by the character identity.

The local density matrix becomes
\begin{equation}
\rho_{\rm cell}=\frac{I_{16}}{16}\otimes\sigma,\qquad
\sigma=\bigoplus_{s=1}^{35}\frac{P_s}{3}I_3.
\label{eq:localfactor}
\end{equation}
Because the original routing-orientation block is proportional to the identity for each $s$, this is the original full density operator written in a new basis. Its entropy decomposes as
\begin{equation}
g_{\rm share}=4\log2+S(\sigma)
=2.77258872224+4.64721130133.
\end{equation}

\subsection{The coherent global state}

Pair each virtual port mode with the mode across its graph link. On a link $e=(L,R)$ choose
\begin{equation}
|\omega_e\rangle=\frac{|10\rangle+|01\rangle}{\sqrt2}.
\end{equation}
The product over links, with a fixed fermion ordering, is
\begin{equation}
|\Omega_\Gamma\rangle=
\prod_{e=(L,R)}
\frac{c^\dagger_{L,e}+c^\dagger_{R,e}}{\sqrt2}|0\rangle,
\qquad
\rho_{\Gamma,*}=|\Omega_\Gamma\rangle\langle\Omega_\Gamma|
\otimes\bigotimes_v\sigma_v.
\end{equation}
Every cell has the marginal \eqref{eq:localfactor}. A boundary cutting $n$ links has bond entropy $n\log2$. Residual register entropy is a volume contribution and is conditioned out of the geometric cut calculation.

The joint coherence is a physical choice. The separable state
\begin{equation}
\rho_{\rm sep}=\tfrac12|10\rangle\langle10|
+\tfrac12|01\rangle\langle01|
\end{equation}
has the same one-side marginal as $|\omega\rangle$, but zero entanglement. The finite local probabilities therefore admit both states.

A local bond parent selects the coherent state:
\begin{equation}
H_e=U(n_L+n_R-1)^2
+t\left[n_L+n_R-2n_Ln_R
-(c_L^\dagger c_R+c_R^\dagger c_L)\right],
\quad U,t>0.
\label{eq:bondparent}
\end{equation}
The zero-energy eigenvector is $|\omega_e\rangle$. The empty and doubly occupied states have energy $U$; the antisymmetric one-particle state has energy $2t$. Equation \eqref{eq:bondparent} provides an explicit interaction realizing the required correlations.

\subsection{Coherent renewal and changing graph frames}

Independent resets of individual cells destroy the intercell bond state. A compatible renewal exports the replaced capacity sector and prepares the entire local bond pattern with the prescribed residual state. For a fixed graph it has a Stinespring isometry.

For a superposition of graphs, a direct export of graph-dependent bond vacua can record which graph was present. Choose preparation unitaries $U_g|\mathrm{ref}\rangle=|\Omega_g\rangle$ and define
\begin{equation}
D=\sum_g|g\rangle\langle g|\otimes U_g,\qquad
J|g\rangle=|g\rangle|\Omega_g\rangle.
\end{equation}
The selected moving-frame rule first applies $D^\dagger$, performs renewal in the common reference frame, evolves the retained geometry, then prepares the new graph frame. For a supplied retained operator $U_{\rm geom}$,
\begin{equation}
U_{\rm phys}=D(U_{\rm geom}\otimes I)D^\dagger,\qquad
U_{\rm phys}J=JU_{\rm geom}.
\label{eq:frame}
\end{equation}
The old record remains exported; the frame change acts before export on the current capacity state.

This construction handles a specified common register organization. When geometric alternatives have different numbers of renewal events or different record spaces, a compatible common dilation must also be given. The later finite orbital Hamiltonian does not settle that enlarged record problem.

\section{Proper ticks, event volume, and the bare gravitational scale}
\label{sec:scale}

\subsection{Causal chains}

Assume a locally Poisson event process with density $\rho=1/v_*$ in invariant four-volume, and identify the physical cadence between interaction endpoints with a longest causal chain. In four-dimensional Minkowski space the causal-diamond volume is
\begin{equation}
V_4(\tau)=\frac{\pi}{24}\tau^4
\end{equation}
in $c=1$ units. By definition of the asymptotic chain constant,
\begin{equation}
L_{\rm chain}\sim m_4(\rho V_4)^{1/4}.
\end{equation}
Matching the asymptotic chain rate to $\tau/\tstar$ gives
\begin{equation}
\boxed{\frac{v_*}{\Lstar^4}=\frac{\pi}{24}m_4^4.}
\label{eq:eventvol}
\end{equation}
The event-placement law and longest-chain cadence are additional premises. They concern spacetime events; the memoryless capacity replacement law concerns marks on those events.

Conditioning the event process on a fixed total number produces a finite-volume binomial law for subregions. For a region of volume fraction $f$ in a total of $N_e$ events, the local event count has variance $N_ef(1-f)$. A small subregion of a much larger system approaches the Poisson law. This provides a consistent distinction between a fixed total event sector and local fluctuations.

The retained numerical value is $m_4=2.35$, which gives
\begin{equation}
v_*/\Lstar^4=3.9921796826723166.
\end{equation}
This value is provisional. The construction does not replace it by an unpublished later estimate or treat a finite-size fit as an exact constant.

\subsection{The paired-cell volume calibration}

Assign one event to a pair of primitive spatial cells over one proper tick. The volume per cell is then
\begin{equation}
V_{\rm cell}=\frac{v_*}{2\Lstar}
\end{equation}
in length units. Select a regular-tetrahedron semiclassical volume reading,
\begin{equation}
V_{\rm cell}=\frac{a^3}{6\sqrt2},\qquad
\Ach=\frac{\sqrt3}{4}a^2.
\end{equation}
Consequently
\begin{align}
\frac{a}{\Lstar}
&=\left(3\sqrt2\,\frac{v_*}{\Lstar^4}\right)^{1/3},\\
\frac{\Ach}{\Lstar^2}
&=\frac{\sqrt3}{4}
\left(3\sqrt2\,\frac{v_*}{\Lstar^4}\right)^{2/3}.
\label{eq:areascale}
\end{align}
The two-cell factor is a physical accounting choice. The regular volume is a semiclassical observable assignment. The continuation exhibits compatible block states, but neither item follows from the integer $1680$ alone.

At $m_4=2.35$,
\begin{equation}
a/\Lstar=2.56812077553,\qquad
\Ach/\Lstar^2=2.85582456165.
\end{equation}
A host edge therefore spans several proper tick lengths. Monte Carlo updates, causal slice labels, edge lengths, and proper ticks retain different meanings.

\subsection{Absolute angular transfer}

One side of a coherent binary bond has reduced state $\rho_B=I_2/2$. Select a physical angular generator $K_B$ by the \emph{absolute} full-turn condition
\begin{equation}
e^{-2\pi K_B}=\rho_B.
\end{equation}
It gives
\begin{equation}
K_B=bI_2,\qquad b=\frac{\log2}{2\pi}.
\label{eq:angular}
\end{equation}
The absolute condition fixes an otherwise free additive scalar in $K_B$. Normalized density matrices alone are invariant under $K_B\mapsto K_B+\lambda I$, because normalization removes the scalar. An entropy calculation by itself would leave this curvature normalization undetermined.

Let the sum of wedge angles around an interior hinge be $2\pi-\delta_h$. For $N_h$ channels,
\begin{align}
\Tr\exp[-(2\pi-\delta_h)K_B]
&=2e^{-b(2\pi-\delta_h)}
=e^{b\delta_h},\\
W_h&=e^{bN_h\delta_h}.
\end{align}
The flat full turn has weight one. A replicated cut with total angle $2\pi n$ has
\begin{equation}
\log\Tr\rho_B^n=(1-n)\log2.
\end{equation}
Thus the selected tensor has both the binary cut entropy and a definite absolute curvature response.

\subsection{The coefficient of the Regge action}

With $A_h=\Ach N_h$, the geometric weight is
\begin{equation}
W_{\rm curv}=\exp\left[\frac{b}{\Ach}\sum_hA_h\delta_h\right].
\end{equation}
The Regge approximation uses $\int\sqrt g\,R=2\sum_hA_h\delta_h$, including the corresponding boundary term when a boundary is present. Comparing $W=e^{-I_E}$ with
\begin{equation}
I_E=-\frac1{16\pi G_0}\int\sqrt g\,R
\end{equation}
gives
\begin{equation}
\boxed{G_0=\frac{\Ach}{8\pi b}=\frac{\Ach}{4\log2}.}
\label{eq:G0}
\end{equation}
No additional Einstein weight is multiplied into the state sum.

Substituting \eqref{eq:eventvol} and \eqref{eq:areascale},
\begin{equation}
\boxed{
\frac{G_0}{\Gstar}
=\frac{\sqrt3}{16\log2}
\left(\frac{\pi\sqrt2}{8}\right)^{2/3}m_4^{8/3}
=0.10551926870131369\,m_4^{8/3}.}
\label{eq:Gm4}
\end{equation}
At the provisional input this gives $G_0/\Gstar=1.03002098319$. Exact equality in this calibration would require
\begin{equation}
m_4=(0.10551926870131369)^{-3/8}\simeq2.32407739.
\end{equation}
That inverse value is a target implied by the equality; it is not a measurement of $m_4$.

Logarithmic differentiation gives
\begin{equation}
\frac{dG_0}{G_0}=\frac83\frac{dm_4}{m_4}.
\end{equation}
A one-percent uncertainty in the asymptotic chain constant produces about a $2.67$-percent uncertainty in this bare coefficient. Numerical roundoff in the finite tensor identities is much smaller than this physical calibration uncertainty.

\subsection{Boundary and bulk consistency}

The cut entropy per physical area is $\log2/\Ach$. With \eqref{eq:G0},
\begin{equation}
\frac{S_{\rm cut}}{A}=\frac1{4G_0}.
\end{equation}
The agreement follows from the common angular tensor. It supplies a consistency identity within the selected microscopic action. A measurement of a renormalized gravitational coupling can differ from its bare coefficient while retaining a consistent renormalized effective action and boundary law.

The main manuscript's C.5-matched channel area is $4\log2\,\Lstar^2$, obtained after selecting $G=\Gstar$. The paired-cell route supplies $\Ach$ independently of that final equality and obtains the nearby bare ratio \eqref{eq:Gm4}. Both routes cannot be described as independent measurements of the same interacting $G_R$.

\section{Primitive geometric cells and connected interfaces}
\label{sec:geometric}

\subsection{Relational geometric data}

A local spatial frame change acts on link transport as
\begin{equation}
g_{cc'}\mapsto h_cg_{cc'}h_{c'}^{-1}.
\end{equation}
Fluxes at a tetrahedron satisfy the Gauss constraint
\begin{equation}
\sum_{f=1}^4\mathbf J_f=0.
\end{equation}
The invariant intertwiner space carries shape information after this constraint is imposed. The capacity statistic $K^2$ continues to fluctuate in its separate factor. Setting it equal to a Gauss Casimir would make it vanish on geometric singlets and destroy its defining Gibbs distribution.

The selected one-face capacity--geometry lock is an equivariant map between two copies of $V_3$. Schur's lemma fixes the form of a nonzero map up to a scalar. Its existence as an interaction is a selected realization condition. For $N$ aligned faces, maximal fusion gives the unique normalized equivariant map
\begin{equation}
F_N:V_3^{\otimes N}\longrightarrow V_{3N},
\qquad
F_N|3,\mathbf n\rangle^{\otimes N}=|3N,\mathbf n\rangle.
\end{equation}
The coarse spin grows through blocking while the primitive spin stays fixed.

The coherent area prescription uses the additive flux on this maximal ray. It can be related to the kinematical $\sqrt{J(J+1)}$ area by the finite observable factor
\begin{equation}
Z_A(N)=\sqrt{\frac{3N}{3N+1}}.
\end{equation}
The state-sum construction uses the resulting assigned coarse areas and normalized shape packets. It does not require simultaneously sharp quantum dihedral-angle eigenstates.

\subsection{The retained paired-cell sector}

Couple two copies of $\mathcal I_3$ through the four-bundle maximal-fusion map
\begin{equation}
T:\mathcal I_3\otimes\mathcal I_3\longrightarrow\mathcal I_6,
\qquad \dim\mathcal I_6=13.
\end{equation}
With $F=(TT^\dagger)^{-1/2}T$, the operator $I_{\rm pair}=F^\dagger$ is an isometry from the thirteen-dimensional coarse space into the forty-nine-dimensional pair.

There is an explicit polynomial parent for this subspace. Define
\begin{align}
X&=\mathbf J_1\cdot\mathbf J_2,\qquad
Y=\mathbf J_1\cdot\mathbf J_3,\qquad Z=-12I-X-Y,\\
u&=(X+4I,Y+4I,Z+4I),\qquad B=\sum_ru_r^2,\qquad Q=i[X,Y],\\
\mathcal D&=\sum_ru_r\otimes u_r,\\
\mathcal E&=\sum_r(u_r^2\otimes u_r+u_r\otimes u_r^2),\qquad
\mathcal T=\sum_ru_r^2\otimes u_r^2.
\end{align}
The dimensionless parent is
\begin{equation}
\begin{split}
\widehat H_{\rm pair}=\frac1{13^4}\big[
&3\,13^4I+13^2(B\otimes I+I\otimes B-8\mathcal D)\\
&+52\mathcal E+B\otimes B-2\mathcal T-144Q\otimes Q
\big].
\end{split}
\label{eq:pairparent}
\end{equation}
The retained exact rational audit gives
\begin{equation}
\widehat H_{\rm pair}\succeq0,\qquad
\ker\widehat H_{\rm pair}=\Ran I_{\rm pair},
\qquad \dim\ker\widehat H_{\rm pair}=13.
\end{equation}
There are $36$ positive eigenvalues and a dimensionless gap about $1.9301241828$. A physical energy scale multiplying the parent remains a separate input.

Within an assigned finite block this parent selects the code. Applying the same parent to overlapping pairs imposes additional correlations. The stored three-cell calculation, included in the reference inputs, gives a nineteen-dimensional common ground space for $H_{12}+H_{23}$ in the $343$-dimensional space, with positive gap about $0.96520$ in parent-energy units. This is a compatibility check for that finite overlap. The present construction uses assigned blocks and separately specifies their geometric gluing.

\subsection{An integer-channel causal regulator}

Euclidean CDT building blocks have spatial squared edges $a^2$ and mixed squared edges $\alpha a^2$. A spatial triangle has area
\begin{equation}
A_{\rm SL}=\frac{\sqrt3}{4}a^2,
\end{equation}
while a triangle with one spatial and two mixed edges has
\begin{equation}
A_{\rm mix}=\frac{\sqrt{4\alpha-1}}4a^2.
\end{equation}
Assigning one channel to a spatial triangle and three to a mixed triangle gives
\begin{equation}
A_{\rm SL}=\Ach,\qquad A_{\rm mix}=3\Ach,
\qquad \boxed{\alpha=7}.
\label{eq:alpha7}
\end{equation}
Both Euclidean four-simplex volumes are nondegenerate:
\begin{equation}
V_{41}=a^4\frac{\sqrt{8\alpha-3}}{96},\qquad
V_{32}=a^4\frac{\sqrt{12\alpha-7}}{96}.
\end{equation}
The stronger threshold is $\alpha>7/12$.

A four-port patch with $V$ cells and $E$ internal links has $B=4V-2E$ external ports. Hence $B$ is even. A mixed tetrahedron with channel counts $(1,n,n,n)$ requires odd $n$ in this patch architecture. The first anisotropic choice after $n=1$ is $n=3$. This parity argument restricts the specified architecture; it does not prove universal uniqueness of $\alpha=7$.

The three coarse tetrahedral interfaces are
\begin{center}
\begin{tabular}{llll}
\toprule Slice split & Channels & Spins & Dimension\\\midrule
$4$--$0$ & $(1,1,1,1)$ & $(3,3,3,3)$ & $7$\\
$3$--$1$ & $(1,3,3,3)$ & $(3,9,9,9)$ & $7$\\
$2$--$2$ & $(3,3,3,3)$ & $(9,9,9,9)$ & $19$\\
\bottomrule
\end{tabular}
\end{center}

\subsection{Explicit invariant tensors}

An orthonormal recoupling basis for $\Inv(V_{j_0}\otimes V_{j_1}\otimes V_{j_2}\otimes V_{j_3})$ is
\begin{equation}
B_k(m_0,m_1,m_2,m_3)
=\sum_M\frac{(-1)^{k-M}}{\sqrt{2k+1}}
C^{kM}_{j_0m_0,j_1m_1}
C^{k,-M}_{j_2m_2,j_3m_3}.
\label{eq:intertwiner}
\end{equation}
The allowed $k$ are the common intermediate spins of the two pairs. Internal spin-three edges use
\begin{equation}
E_{mm'}=(-1)^{3-m}\delta_{m,-m'}.
\end{equation}
The normalized maximal fusion of three spin-three legs into spin nine is particularly simple in occupation coordinates $r_i=m_i+3$:
\begin{equation}
(F_3)_{R,r_1r_2r_3}
=\delta_{R,r_1+r_2+r_3}
\left[
\frac{\binom6{r_1}\binom6{r_2}\binom6{r_3}}
{\binom{18}R}
\right]^{1/2}.
\label{eq:fusion3}
\end{equation}
The binomial convolution identity gives $F_3F_3^\dagger=I_{19}$. Equations \eqref{eq:intertwiner} and \eqref{eq:fusion3} fix phases and normalization for independent implementations.

\subsection{The four-cell chain}

Take four primitive cells with ports
\begin{equation}
(a,b,c,x),\quad(x,d,e,y),\quad(y,f,g,z),\quad(z,h,i,j).
\end{equation}
Contract repeated internal labels with $E$. Bundle external ports into
\begin{equation}
a;\qquad(b,d,f);\qquad(c,e,h);\qquad(g,i,j).
\end{equation}
Apply \eqref{eq:fusion3} to the three bundles. The resulting map is
\begin{equation}
R_{31}:\mathcal I_3^{\otimes4}
\longrightarrow\Inv(V_3\otimes V_9^{\otimes3}).
\end{equation}
It has rank seven. Restricting consecutive cell pairs to the retained paired sectors gives a $7$-by-$169$ map with the same rank.

The calculation exploits the maximal-spin selection rule at the tail: both last-cell pair channels must be six in the unrestricted recoupling coordinates. An independent residual check of the forbidden channels is about $5.01\times10^{-17}$. This both reduces the contraction cost and checks the implemented index grouping.

\subsection{The six-cell cycle}

Take six cells around a cycle, each with two internal and two external ports. Assign the external port pairs to coarse faces
\begin{equation}
(0,1),(0,2),(1,2),(1,3),(2,3),(0,3).
\end{equation}
Each face receives three primitive facets. Maximal fusion gives four spin-nine faces. Retaining three consecutive paired sectors produces
\begin{equation}
R_{22}:\C^{13}\otimes\C^{13}\otimes\C^{13}
\longrightarrow\Inv(V_9^{\otimes4}),
\end{equation}
a $19$-by-$2197$ map of full row rank.

The implementation evaluates coherent-state polynomials instead of forming a tensor with twelve free magnetic indices. Eighty evaluations determine the map in the orthonormal target basis; sixteen fresh evaluations test it. The evaluation matrix has condition number about $4.11$, and the largest held-out relative residual is $3.29\times10^{-15}$. These are numerical rank and contraction certificates with ample singular-value separation, rather than exact symbolic rank proofs.

The cycle is a connected spin network with grouped boundary facets. A literal tetrahedral three-ball imposes additional geometric incidence requirements. The tested six-tetrahedron fan and an eight-tetrahedron refinement have rank thirteen for the required nineteen-dimensional boundary space. Their failure limits those particular fine-manifold realizations. The full-rank abstract cycle remains the retained construction.

\subsection{Normalized interfaces and bounded parents}

\begin{proposition}[Normalized interface]
Let $R_e:\mathcal H_{\rm patch}\to\mathcal H_e$ have full row rank. Then
\begin{equation}
F_e=(R_eR_e^\dagger)^{-1/2}R_e,\qquad I_e=F_e^\dagger
\label{eq:polar}
\end{equation}
satisfy $F_eI_e=I_{\mathcal H_e}$ and $I_e^\dagger I_e=I_{\mathcal H_e}$.
\end{proposition}
\begin{proof}
The Gram matrix $R_eR_e^\dagger$ is strictly positive. Multiplying \eqref{eq:polar} by its adjoint gives
$F_eF_e^\dagger=(R_eR_e^\dagger)^{-1/2}
(R_eR_e^\dagger)(R_eR_e^\dagger)^{-1/2}=I$.
\end{proof}

Inside the paired subspace, add
\begin{equation}
H_{\rm code}=\Delta_{\rm p}(P_{\rm pair}-P_{\rm code}),
\qquad P_{\rm code}=I_eI_e^\dagger,\quad \Delta_{\rm p}>0.
\end{equation}
Because $P_{\rm code}\le P_{\rm pair}$, this term is positive. Together with the disjoint pair parents, it leaves the desired interface code. The interaction has bounded support, four or six cells in the displayed examples.

The Gram condition numbers are about $9.68$ and $25.81$. Isometry errors are below $2.1\times10^{-14}$. The metric meaning of a code observable is given by transporting the coarse operator through $I_e$. Equality to a sum of unchanged fine-cell volume operators is a further requirement, absent from this construction.

\section{The microscopic state sum and its exact CDT correspondence}
\label{sec:statesum}

\subsection{Causal histories and the coarse measure}

A history is a causal triangulation built from $(4,1)$, $(3,2)$, and their time-reversed four-simplices between adjacent spatial slices. The spatial topology and boundary markings are fixed as part of the ensemble. Standard CDT supplies the causal simplicial history class, its Wick-rotation prescription, and transfer-matrix framework \cite{cdt}. The present construction supplies microscopic tensors whose Euclidean contraction reproduces the chosen weight on that class.

For an unmarked history, write the coarse partition function as
\begin{equation}
Z_{\rm CDT}[J]=\sum_{T\in\mathfrak T}
\frac{1}{C_T}\,e^{-S_E[T]}W_J[T].
\label{eq:Zcdt}
\end{equation}
Here $C_T$ is the appropriate relative automorphism factor, $J$ specifies the source and preparation, and any volume restriction is included in $\mathfrak T$ or in the weight. A labelled enumeration requires its own counting convention; multiplying an additional label degeneracy into an already normalized unlabelled sum changes the measure.

\subsection{Local simplex tensors}

For each coarse tetrahedron $e$, use its CDT edge lengths to construct outward area normals and a normalized coherent intertwiner $u_e$. Its microscopic packet is $\widetilde u_e=I_eu_e$. A spatial tetrahedron uses the one-cell identity map. The mixed tetrahedra use the connected maps of Section \ref{sec:geometric}.

Choose the four-simplex tensor to be the tensor product of its boundary packets multiplied by
\begin{equation}
w_\sigma=\exp\left[-b\sum_{h\subset\sigma}N_h\theta_{h\sigma}
-\lambda V_\sigma\right].
\end{equation}
The packet on the other side of a shared tetrahedron is its adjoint in the same ordering. Supply the factor $e^{2\pi bN_h}$ for an interior hinge and $e^{\pi bN_h}$ for a boundary hinge.

\begin{theorem}[Exact Euclidean weight]
\label{thm:weight}
For every allowed coarse triangulation, contraction of these normalized microscopic tensors gives
\begin{equation}
W_{\rm micro}(T)=
\exp\left[
b\sum_{h\in{\rm int}}N_h\delta_h
+b\sum_{h\in\partial}N_h\psi_h
-\lambda\sum_\sigma V_\sigma
\right].
\label{eq:weightidentity}
\end{equation}
Here $\delta_h=2\pi-\sum_{\sigma\supset h}\theta_{h\sigma}$ and
$\psi_h=\pi-\sum_{\sigma\supset h}\theta_{h\sigma}$.
\end{theorem}
\begin{proof}
Each internal tetrahedron contributes
$\widetilde u_e^\dagger\widetilde u_e
=u_e^\dagger I_e^\dagger I_eu_e=1$.
The remaining simplex angle terms collect independently at each hinge. Adding its full-turn or boundary half-turn factor produces the corresponding deficit or exterior angle. Volume terms add over simplices. No step depends on the number of simplices.
\end{proof}

The rank-one packets fix the shapes to the selected CDT building blocks while allowing connectivity to vary. The theorem is an explicit local tensor realization of that model. It does not identify these chosen tensors with a different spin-foam vertex.

At $A_h=\Ach N_h$, \eqref{eq:weightidentity} is precisely the Regge curvature weight with $G_0$ in \eqref{eq:G0}. The local scalar source check multiplies the same action $W_J[T]$ into both sides before contraction. Thus the equality persists pointwise in the source field. Integrating a dynamical field requires retaining its determinant in both measures.

\subsection{Why the full normalization matters}

Raw packets $R_e^\dagger u_e$ have squared norms about
\begin{equation}
\|R_{31}^\dagger u_{31}\|^2=5.3131575762\times10^{-6},\qquad
\|R_{22}^\dagger u_{22}\|^2=2.1629219768\times10^{-10}.
\end{equation}
These factors multiply at internal tetrahedra. Their product depends on the composition and shape of the triangulation. It cannot be removed by one common normalization of $Z$.

The Gram correction in \eqref{eq:polar} removes the entire unwanted matrix factor. This is the operative measure repair. It determines how a selected coarse CDT model is realized in the microscopic patch space. It neither alters prior campaign data nor establishes that a previously simulated raw measure already had this normalization.

\subsection{Boundary sewing and observables}

Let $I_g$ encode an entire boundary state with consistent symmetry factors. If the coarse transfer is $K[g_2,g_1;J]$, the microscopic transfer on its prepared code is
\begin{equation}
K_{\rm micro}[J]=I K_{\rm CDT}[J]I^\dagger.
\label{eq:transferlift}
\end{equation}
For two intervals, contraction over an orthonormal encoded intermediate boundary gives
\begin{equation}
IK_2I^\dagger IK_1I^\dagger=I(K_2K_1)I^\dagger.
\end{equation}
The normalized boundary metric is therefore part of the sewing identity.

\begin{theorem}[Observable and phase inheritance]
\label{thm:inherit}
Suppose the two theories use the same causal history set, topology, boundary preparation, symmetry factors, volume treatment, and source functional. Suppose their microscopic contraction satisfies \eqref{eq:weightidentity} and their boundary maps satisfy \eqref{eq:transferlift}. Then every normalized source derivative and every observable transported through $I$ agrees in the two descriptions. Any limit of those observables, including a de Sitter profile or a gravitational pole residue, agrees whenever the corresponding limit exists for that same coarse ensemble.
\end{theorem}
\begin{proof}
Substitution in the history sum gives equal partition functions with equal source dependence, up to a common source-independent normalization. Such a factor cancels in normalized expectations. Differentiating the identical generating functions gives identical connected correlators. The boundary statement follows from isometric sewing. Equal regulated observables have equal limits along the same sequence of ensembles whenever that sequence converges.
\end{proof}

The theorem gives a precise route from microscopic construction to coarse physics. Its antecedent includes the actual physical ensemble. Proving the equality of weights does not supply a phase measurement at a new coupling point.

\subsection{The finite causal check}

The test family uses the spatial boundary of a four-simplex, with five tetrahedra, and its product with one interval. For each of the $120$ vertex orderings, staircase triangulation produces a causal slab. Legal $2\to4$ moves add $1800$ distinct neighbors, giving $1920$ marked histories.

A proposed move replaces $A*\partial B$ by $\partial A*B$, with $|A|=4$ and $|B|=2$. The new edge $B$ must be absent; the old star must have the required two simplices; new simplices must span adjacent slices; boundary tetrahedra must retain the prescribed boundary. These tests check the local ball replacement and causal slicing. The finite family is a stated restriction of the full history sum.

The primitive tensor log weight agrees with direct Regge evaluation to $1.46\times10^{-13}$. The external scalar action comparison reaches $1.43\times10^{-13}$. Two-slab curvature sewing is zero at displayed precision, the scalar quadratic-kernel sewing error is $1.07\times10^{-14}$, and the sewn source covariance error is $2.22\times10^{-16}$. These checks test the identity on actual nontrivial causal complexes.

\section{From the selected curvature action to Einstein response}
\label{sec:einstein}

\subsection{Metric variation and collective areas}

For a length-metric Regge complex, the Schläfli identity
\begin{equation}
\sum_{h\subset\sigma}A_h\,d\theta_{h\sigma}=0
\end{equation}
removes the angle-variation terms in the curvature action. Therefore
\begin{equation}
\delta S_E=-\frac1{8\pi G_0}\sum_h\delta_h\,\delta A_h
+\lambda\sum_\sigma\delta V_\sigma+\delta S_{\partial}.
\end{equation}
On a controlled refinement of smooth metrics, this is the discretization of the Einstein variational problem. In the fixed-edge CDT model, geometry changes through connectivity. Passing from its sum to a smooth metric effective action is a dynamical coarse-graining question.

Six independent directional areas can reconstruct a positive spatial metric. Let $C$ be the cofactor matrix of $g$:
\begin{equation}
C=(\det g)g^{-1},\qquad
A(n)^2=n^TCn.
\end{equation}
Use normals $e_x,e_y,e_z,(e_x+e_y)/\sqrt2,(e_x+e_z)/\sqrt2,(e_y+e_z)/\sqrt2$. The first three areas give the diagonal entries of $C$; for example
\begin{equation}
C_{xy}=A((e_x+e_y)/\sqrt2)^2-\frac{C_{xx}+C_{yy}}2.
\end{equation}
On $C>0$,
\begin{equation}
g=\sqrt{\det C}\,C^{-1},\qquad
\sqrt{\det g}=(\det C)^{1/4}.
\end{equation}
This provides a metric dictionary for collective channels. The dictionary must refer to physical coarse cuts when used in a geometry ensemble.

\subsection{The quadratic metric action}

Expand the selected Lorentzian Einstein action about flat space:
$g_{\mu\nu}=\eta_{\mu\nu}+h_{\mu\nu}$, with signature $(-,+,+,+)$.
Gauge transformations are
\begin{equation}
\delta h_{\mu\nu}=\partial_\mu\xi_\nu+\partial_\nu\xi_\mu.
\end{equation}
The Hessian annihilates these directions. After the lapse and shift constraints and residual gauge equivalence are treated, an on-shell null momentum has two transverse-traceless polarizations. In the usual metric perturbation normalization their kinetic action is
\begin{equation}
S_{\rm TT}^{(2)}
=\frac1{64\pi G_0}\int d^4x\,
\left[(\partial_t h_{ij}^{\rm TT})^2
-(\partial_kh_{ij}^{\rm TT})^2\right].
\end{equation}
The positivity of these physical modes coexists with the familiar indefinite scalar direction before constraint reduction.

The saved Fourier calculation constructs the length-Regge Hessian from twenty-four Kuhn simplices per unit hypercube. It checks flat deficits, gauge null vectors, tensor orientations, and the small-momentum Fierz--Pauli limit. Separate off-shell and on-shell mode counts are retained: the Euclidean off-shell physical tensor block contains more components than the two propagating Lorentzian helicities.

\subsection{The static source and Newton's law}

Use one metric source coupling,
\begin{equation}
S_{\rm int}=\frac12\int d^4x\,h_{\mu\nu}T^{\mu\nu}.
\end{equation}
Variation of the full parent gives
\begin{equation}
G_{\mu\nu}+\Lambda g_{\mu\nu}=8\pi G_0T_{\mu\nu}.
\label{eq:einstein}
\end{equation}
For a weak, nonrelativistic, static source with negligible anisotropic stress,
\begin{equation}
ds^2=-(1+2\Phi)dt^2+(1-2\Psi)d\mathbf x^2,\qquad
\Phi=\Psi,
\end{equation}
and the $00$ constraint gives
\begin{equation}
\nabla^2\Phi=4\pi G_0\rho.
\end{equation}
For a point mass, use $\nabla^2(1/r)=-4\pi\delta^{(3)}(\mathbf x)$:
\begin{equation}
\Phi(r)=-\frac{G_0M}{r},\qquad
\mathbf a=-\nabla\Phi=-\frac{G_0M}{r^2}\widehat{\mathbf r}.
\label{eq:newton}
\end{equation}
The same coefficient multiplies the metric kinetic term and the force.

In the Euclidean conserved-source calculation, write $g_{44}=1+h_{44}^{E}$ and use the insertion $\delta S_m^E=\tfrac12\int \rho h_{44}^{E}$. The static response is
\begin{equation}
h_{44}^{E}(\mathbf k)=-\frac{8\pi G_0\rho(\mathbf k)}{|\mathbf k|^2},
\qquad \Phi=\frac12h_{44}^{E}.
\end{equation}
The Lorentzian component $h_{00}^{L}=-2\Phi$ has the opposite sign. The saved small-momentum source calculation uses the Euclidean convention and approaches the predicted residue as the lattice momentum decreases. Its convergence checks validate the selected metric Hessian and source vertex. They do not measure the interacting CDT propagator.

Sampling the two wave polarizations alone gives no static density response. The static potential is carried by the constrained sector. Eliminating constraints must therefore preserve the source before an effective response is inferred.

\subsection{Capacity as a metric observable}

For a static metric $ds^2=-N^2dt^2+h_{ij}dx^idx^j$ with $\Lambda=0$, the Einstein equations imply
\begin{equation}
D^2N=4\pi G_0N(E+S),
\end{equation}
where $E$ is the energy density measured by the normal observer and $S$ the spatial stress trace. Define $q=N^2$ and $s=S_\infty(1-q)$. Then
\begin{equation}
D^2q-\frac{(Dq)^2}{2q}=8\pi G_0q(E+S),
\end{equation}
or
\begin{equation}
D^2s+\frac{(Ds)^2}{2S_\infty(1-s/S_\infty)}
=-8\pi G_0S_\infty(1-s/S_\infty)(E+S).
\label{eq:capacitymetric}
\end{equation}
In the weak-field nonrelativistic regime, $s=-2S_\infty\Phi$ and the equation reduces to the manuscript's Poisson bridge.

The lapse depends on clock and boundary normalization. Inside a supported spherical shell, local curvature can vanish while a clock is redshifted relative to infinity. A curvature scalar alone therefore cannot supply this boundary-relative clock deficit. Equation \eqref{eq:capacitymetric} is one constrained metric equation; the spatial equations and matter equilibrium remain part of the same solution.

\section{De Sitter geometry and the phase condition}
\label{sec:desitter}

\subsection{The classical de Sitter solution of the selected action}

For a homogeneous closed spatial metric,
\begin{equation}
ds^2=-dt^2+a_{\rm cos}(t)^2d\Omega_3^2,
\end{equation}
the vacuum Einstein equations with $\Lambda>0$ give
\begin{equation}
\left(\frac{\dot a_{\rm cos}}{a_{\rm cos}}\right)^2
+\frac1{a_{\rm cos}^2}=\frac{\Lambda}{3}.
\end{equation}
Writing $\ell=\sqrt{3/\Lambda}$, the solution is
\begin{equation}
a_{\rm cos}(t)=\ell\cosh(t/\ell).
\end{equation}
Its regular Euclidean section has
\begin{equation}
a_E(\tau)=\ell\cos(\tau/\ell),\qquad
V_3(\tau)=2\pi^2\ell^3\cos^3(\tau/\ell),
\label{eq:dsprofile}
\end{equation}
with total volume $V_4=8\pi^2\ell^4/3$. These equations follow from the selected continuum action and boundary conditions.

A fixed-volume saddle can produce the same Euclidean geometry through a Lagrange multiplier. The multiplier enforcing volume is different from a claim that the microscopic physical cosmological term has already been derived. Appendix O supplies a separate fixed-event-measure premise and source-clock release construction.

\subsection{What CDT evidence establishes}

CDT simulations have exhibited an extended regime with a macroscopic de Sitter volume profile and reconstructed an effective volume action from fluctuations \cite{ds}. This is evidence for those sampled ensembles and observables. The tensor equality in Theorem \ref{thm:weight} makes the microscopic realization inherit such evidence when its ensemble, parameters, and physical dictionary coincide with the tested CDT model.

For the present $\alpha=7$ branch, the displayed zero-volume-term couplings are
\begin{equation}
(\kappa_0,\Delta_{\rm CDT},\kappa_4)
=(2.20205311516,\;0.326138328034,\;-0.304774619960).
\end{equation}
No extended-phase measurement at this point belongs to the current finite construction. The curvature-only $\kappa_4$ is also distinct from a critical or volume-controlled value used in a simulation. Equation \eqref{eq:dsprofile} is a classical target, and Theorem \ref{thm:inherit} is a conditional transfer of ensemble results.

\subsection{Volume stiffness and physical Newton response}

An effective volume action is often fitted in the form
\begin{equation}
S_{\rm vol}\simeq k_1\sum_t
\frac{(n_{t+1}-n_t)^2}{n_{t+1}+n_t}
+\sum_tU(n_t),
\end{equation}
with additional factors depending on the stated volume and time conventions. Its Hessian predicts the covariance after volume constraints and alignment modes are handled. Off-diagonal precision-matrix entries can isolate the nearest-neighbor kinetic term if the assumed local form is valid. Finite samples, projection, volume fixing, and nonlocal effective terms require separate checks.

The earlier blob conversion eliminated several conventions and gave
\begin{equation}
\frac{G_{\rm vol}}{\Lstar^2}
=\frac{\widetilde s_0^2}{3\sqrt6\,k_1}
\sqrt{n_e\,\frac{v_*}{\Lstar^4}},
\end{equation}
where $n_e$ is the physical event count per simplex in that dictionary. This formula does not determine $n_e$ or prove that the fitted homogeneous mode has the same residue as a physical mass-source response. In particular, event counting and simplex counting need a physical map.

At finite physical cutoff, the immediate target is an extended regime with a scale window $a\ll r\ll L_{\rm box}$, appropriate metric tensor behavior, and a source response. Removing the cutoff through a critical surface is a stronger possible goal. This distinction agrees with Appendix J.9. A de Sitter average profile alone does not determine all local propagating modes or their coupling to matter.

\subsection{Registered connected-host results}
\label{sec:alpha7results}

The connected $\alpha=7$ construction fixes the dimensionless curvature
normalization and the two curvature couplings,
\[
 \frac{G_0}{a^2}=\frac{\sqrt3}{16\ln2}=0.1561763194,\qquad
 (\kappa_0,\Delta)=(2.20205311516,0.326138328034).
\]
The associated geometry weight is sampled with the capacity gate off.
The cosmological fugacity and volume feedback remain ensemble controls.
The registration of 17 October 2026 and its amendment were frozen before
sampling at these couplings. The amendment made $N_{41}=32{,}000$ the
primary volume and retained the 16k and 64k runs as reported-only checks.
The target counts $N_{41}$ simplices; the total number of four-simplices is
$N_4=N_{41}+N_{32}$.

\paragraph{Registered observables.}
Q1 tests whether the de~Sitter-like blob is independent of its starting
configuration. Q2 compares the Newton coefficient inferred from the
homogeneous volume stiffness with the boundary coefficient $G_0$:
\[
 R\equiv\frac{G_{\rm vol}}{G_0}
 \simeq 0.076503\,\frac{\widetilde s_0^{4/3}}{k_1}.
\]
The width $\widetilde s_0$ and bias-corrected local stiffness $k_1$ are
measured; their conversion is fixed by the registered dictionary. The
ratio is independent of the provisional chain constant $m_4$. No
renormalized stiffness or source response was tuned to make $R=1$.
The written prior was $R\simeq1.4$, with disagreement considered more likely.

The final window contains 200,000 sweeps per run. Blob eligibility requires
a relative profile-fit error at most $0.05$, width $2.5\leq w\leq K/8$,
and every registered half-window stationarity statistic to satisfy
$|z|\leq3$. Here $K=80$ is the number of time slices. Q2 uses eligible
32k runs only. Its registered agreement rule places the central geometric
mean in $0.80\leq\overline R\leq1.25$; it does not require the entire
uncertainty interval to lie inside that band.

\begin{center}
\small
\begin{tabular}{lrrrll}
\hline
Run & $N_{41}$ target & $R$ & $\sigma_{\ln R}$ & Blob eligibility & Role\\
\hline
A3 & 32,000 & 1.198 & 0.059 & Drifting & Primary pool\\
A4 & 32,000 & 1.064 & 0.089 & Pass & Primary pool\\
A6 & 32,000 & 1.027 & 0.062 & Drifting & Primary pool\\
A1 & 16,000 & 1.164 & 0.061 & Drifting & Reported only\\
A2 & 16,000 & 1.060 & 0.065 & Pass & Reported only\\
A5 & 64,000 & 0.893 & 0.054 & Pass & Reported only\\
\hline
\end{tabular}
\end{center}

\paragraph{Normalization result.}
A4 is the sole eligible primary run. It gives
\[
 \overline R=1.064,\qquad \sigma_{\ln\overline R}=0.089,
 \qquad \text{Q2: AGREE}.
\]
The approximate two-standard-error interval is $0.89$--$1.27$.
Its relative profile-fit error is $0.013$. A2 and A5 also have stationary
de~Sitter-like profiles, with fit errors $0.015$ and $0.012$, and their
ratios lie inside the same agreement band. They support the volume
comparison without changing the primary verdict. Each eligible volume
has one stationary run. The change from A4 to A5 is $1.7$ standard errors
in $\ln R$; a large-volume limit has not been estimated from these data.
The specified microscopic history measure thus has a finite-volume
stiffness compatible with its own boundary normalization under the
registered conversion.

\paragraph{Phase and estimator checks.}
Q1 remains \textsf{undetermined}. A3 fails stationarity through $\xi$
at $|z|=3.3$, and A6 through $\widetilde s_0$ at $|z|=3.3$.
A3 and A4 also differ in largest-slice fraction by $6.1$ standard errors.
The existence of stationary blobs therefore coexists with an unresolved
test of start independence. A post hoc inverse-covariance calculation,
after removing the active small-sample correction floor, gives
$R=1.179,1.022,0.973$ for A2, A4 and A5. These are exploratory estimator
checks; the local-stiffness values above determine the registered result.

\paragraph{Matter-response status.}
A separately registered curvature-trace test used two new streams at
fixed $\kappa_4$, with 64 saved geometries and two source marks per
geometry in each stream. Its dimensionless diagnostic was
$\gamma_{\rm trace}=1.82\pm1.04$, with approximate 95\% interval
$[-0.65,4.28]$, against the target $G_0/a^2=0.1562$.
Both streams also exceeded the registered maximum blob width of 10,
with widths $11.17$ and $10.61$. The verdict is
\textsf{undetermined}; no matter Newton coefficient was reported.
This separate fixed-fugacity gate failure leaves Q2 unchanged.

Subsequent quenched-scalar measurements on those 128 saved geometries
are development analyses. They use both saved sources and six bare
masses. The inferred binding slopes depend on distance range, binning
and scalar discretization, while fixed nongravitating controls also
produce apparent short-range slopes. A conditional-variance diagnostic
does not restore a stable binding interpretation. These measurements
supply no accepted binding energy or gravitational null result.
Further matter-response production is deferred pending a calibrated
finite-geometry observable and an independent power assessment.

The completed evidence concerns the microscopic realization, stationary
finite-volume profiles and the registered homogeneous stiffness ratio.
Equality with a matter-source Newton coefficient, a common matter cone,
and the interacting massless transverse--traceless sector remain open.
The illustrative $G_0/G_*=1.030$ at $m_4=2.35$ retains its separate
event-volume and chain-clock premises.


\subsection{The force observable}

For mass-independent held curves with dimensionless Euclidean lengths $\ell_1,\ell_2$, define
\begin{equation}
Z(m_1,m_2)=\langle e^{-m_1\ell_1-m_2\ell_2}\rangle_0,\qquad
\mathcal R=\frac{Z(m_1,m_2)Z(0,0)}{Z(m_1,0)Z(0,m_2)}.
\end{equation}
At zero mass,
\begin{equation}
\partial_{m_1}\partial_{m_2}\log\mathcal R
=\Cov_0(\ell_1,\ell_2).
\label{eq:cov}
\end{equation}
If a stable large-duration interaction energy exists,
\begin{equation}
V_{\rm int}(r)=-\lim_{T\to\infty}\frac1T\log\mathcal R,
\qquad
\frac{V_{\rm int}(r)}{m_1m_2}
=-\frac{\Cov_0(\ell_1,\ell_2)}{T}+O(m_i).
\end{equation}
In an Einstein regime this tends to $-G_R/r$ in compatible physical units.

The preparation must hold relational separation in the fluctuating metric. Fixing each curve's proper length separately removes the covariance in \eqref{eq:cov}. Discarding geometries where source construction fails reweights the ensemble. All requested geometries must therefore have controlled source coverage, or the whole-ensemble estimate must remain unreported.

A small, drifting toroidal sample with sub-edge separations is a finite diagnostic. Its completion would test source machinery at those scales. An infrared $G_R$ requires a resolved distance range, duration stability, mixing, and the specified ensemble. The new normalized $\alpha=7$ construction and older $\alpha=11.312$ runs retain distinct definitions.

\section{Native proper-time sources and coherent renewal}
\label{sec:clock}

\subsection{Activity along a causal curve}

The retained rest generator attaches to a future timelike curve by the selected H.9a proper-time prescription:
\begin{equation}
U_e[\gamma]=\exp[-i\eps P_e\,\tau_\gamma/\tstar].
\end{equation}
For consecutive segments, $\tau_{\gamma_2\circ\gamma_1}=\tau_{\gamma_2}+\tau_{\gamma_1}$, so their transfers compose. The same $\eps$ fixes rest energy and the proper-time phase; no additional gravitational charge enters.

An energy eigenstate acquires a phase. Interference of two paths reads their relative proper-time phase. This construction does not require a superposition of charged matter and vacuum or an ideal unbounded clock register.

For the $\alpha=7$ causal slabs, two fixed initial positions and tangents launch future normal geodesics. Intrinsic simplex metrics determine propagation and interface transport. Every one of the $1920$ histories is retained. The first duration ranges from about $6.53686$ to $6.96412$ ticks, and the second from $6.53754$ to $6.98966$. The native logarithmic rest-transfer sewing error is $1.72\times10^{-16}$; the local metric transport error is $5.89\times10^{-15}$.

The local Wick check holds a curve segment fixed while continuing its metric quadratic form. Its endpoint is the principal-root relation between Euclidean and Lorentzian segment lengths. This verifies the local formula used on those segments. A reflection-positivity theorem for the entire marked path-selection rule would also have to control the boundary preparation and the sum over paths.

\subsection{A source that is identically insensitive}

A path following one fixed-length timelike CDT edge per slice has length
\begin{equation}
\ell=Ka\sqrt\alpha.
\end{equation}
If its mass dependence is only $e^{-m\ell}$, this factor is constant across the geometry ensemble and cancels from the connected two-source ratio. The mixed mass derivative vanishes exactly at every system size.

The intrinsic free curves above have geometry-dependent lengths. Their selected source partition function gives
\begin{equation}
\Cov_0(\ell_1,\ell_2)=-0.00171521593106.
\end{equation}
At the native gap,
\begin{equation}
\log\mathcal R(\eps,\eps)
=-3.00472647\times10^{-48}+O(\eps^3).
\end{equation}
The finite preparation is nonstatic, so this sign describes a free-line covariance. A static Newton force requires the additional held-source limit of Section \ref{sec:desitter}.

\subsection{Branch-dependent resets}

Let the system register and each fresh bath slot initially have state $\rho_*$. One retained particle branch triggers $n$ sequential SWAP resets and another triggers $m$. Composing the two register permutations leaves an unmatched cycle of length $|n-m|+1$. The permutation trace identity gives the coherence multiplier
\begin{equation}
\gamma_{nm}=\Tr\rho_*^{|n-m|+1}.
\label{eq:resetoverlap}
\end{equation}
For one extra reset, $\gamma=0.0006036303$ in the full $1680$-state capacity distribution, or $0.00965808$ for the $105$-state residual distribution. The stationary marginal can be correct while the exported records strongly suppress interference.

The free-sector choice renews the background cells according to their own local event schedule, independently of the particle's unmeasured occupancy. Position and the internal mismatch fiber remain retained:
\begin{equation}
\mathcal H=\mathcal H_{\rm ren}\otimes\mathcal H_X
\otimes(\C|0\rangle\oplus Q_{\rm mis}).
\end{equation}
For a fixed geometry the renewal channel is
\begin{equation}
\mathcal E=\mathcal E_{\rm background}\otimes\operatorname{id}_{X,Q},
\end{equation}
and, even for correlated initial states,
\begin{equation}
\Tr_{\rm background}\mathcal E(\rho)
=\Tr_{\rm background}\rho.
\end{equation}
The finite check reaches $5.73\times10^{-17}$. Two orders of spacelike-disjoint background resets agree to $3.63\times10^{-17}$. These equalities check the specified local channel. The gravitational deformation algebra and coherence between different geometric event counts remain separate interfaces.

\subsection{A sufficient operator clock}

\begin{theorem}[Positive clock from a half-transfer]
\label{thm:clock}
Let $A$ be a finite half-transfer and let
$0<N_{\min}I\preceq N\preceq N_{\max}I$ on its intermediate space. On the fixed support of
\begin{equation}
K(m)=A^\dagger e^{-mN}A,\qquad
H(m)=-\frac1{\Delta t}\log K(m),
\end{equation}
the clock rate obeys
\begin{equation}
\frac{N_{\min}}{\Delta t}I
\preceq\partial_mH
\preceq\frac{N_{\max}}{\Delta t}I.
\label{eq:clockbound}
\end{equation}
\end{theorem}
\begin{proof}
Set $C=-\partial_mK=A^\dagger Ne^{-mN}A$. Functional calculus for $N$ gives $N_{\min}K\preceq C\preceq N_{\max}K$. The Fréchet derivative of the logarithm is
\begin{equation}
\mathcal D\log_K(C)=\int_0^\infty
(K+sI)^{-1}C(K+sI)^{-1}\,ds.
\end{equation}
This linear map preserves operator order and sends $K$ to $I$. Since $\partial_mH=\mathcal D\log_K(C)/\Delta t$, the claimed bounds follow.
\end{proof}

For $U_t=e^{-itH(m)}$, differentiation of the exponential yields
\begin{equation}
\mathcal T(t)=iU_t^\dagger\partial_mU_t
=\int_0^tU_s^\dagger(\partial_mH)U_s\,ds.
\label{eq:integratedclock}
\end{equation}
It is positive for $t\ge0$ and obeys
\begin{equation}
\mathcal T(t+s)=\mathcal T(t)+U_t^\dagger\mathcal T(s)U_t.
\end{equation}
A supplied constant duration gives the corresponding constant rest-generator shift. A seven-state matrix example checks the formulas through the actual $169$-by-seven geometric encoding.

The operator structure is essential. A retained three-state counterexample has a positive definite transfer with positive assigned entry lengths but a negative eigenvalue of $-\partial_m\log K$. Entrywise positivity therefore does not supply \eqref{eq:clockbound}.

The inverse-lapse clock representation requires a positive lapse on the states where it is used. At zero lapse, that inversion can fail while a self-adjoint reference-time Hamiltonian remains unitary. A turning point of spatial volume alone need not stop the clock.

\subsection{History compression}

A scalar duration mixture satisfies
\begin{equation}
-\partial_m\log\sum_Tp_Te^{-m\ell_T}
=\langle\ell\rangle_m.
\end{equation}
Writing the sum as $Z(m)$ gives $\partial_m^2\log Z=\operatorname{Var}_m(\ell)$, or $\partial_m\langle\ell\rangle_m=-\operatorname{Var}_m(\ell)$. Replacing the distribution by one mean length preserves the first derivative at one point but loses the variance. For the first tested curve the mean is $6.72557455$ ticks and the variance $0.00718245784$.

An enlarged duration register can represent the mixture. Its labels must match the physical boundary organization. If $A_{yx}=\sum_h a_h$ is a coherent half-transfer, assigning an orthogonal intermediate state to every history changes
\begin{equation}
\sum_y\left|\sum_{h:x\to y}a_h\right|^2
\quad\text{into}\quad
\sum_{y,h:x\to y}|a_h|^2.
\end{equation}
The removed cross terms change the kernel. The finite orbital construction below supplies its own common Hilbert space and Hamiltonian explicitly.

\section{One finite geometry--orbital--clock operator}
\label{sec:joint}

\subsection{Cochains and the spatial generator}

Let $C^p(g)$ be the $p$-cochains of a spatial complex and $d_p$ its oriented incidence coboundary. Choose positive metric matrices $M_p$. The metric adjoint is
\begin{equation}
d_p^*=M_p^{-1}d_p^\dagger M_{p+1}.
\end{equation}
The Dirac--Kähler operator $D=d+d^*$ is self-adjoint in this metric. If $M_p=L_pL_p^\dagger$ is a Cholesky factorization, orthonormal coordinates are $y_p=L_p^\dagger x_p$, and adjacent-degree blocks become
\begin{equation}
B_p=L_{p+1}^\dagger d_pL_p^{-\dagger}.
\end{equation}
The full matrix is block tridiagonal in degree with $B_p$ and $B_p^\dagger$ off the diagonal. The identity $d_{p+1}d_p=0$ is checked directly from incidence.

Select the positive-energy orbital operator
\begin{equation}
\Omega_g(m)=\sqrt{D_g^2+m^2I}.
\label{eq:omega}
\end{equation}
It acts on $\mathcal H_X$ and identically on the retained nine-state fiber in this free example. A harmonic cochain has energy $m$.

The positive square root retains components from both signs of the original Dirac--Kähler spectrum. A physical single-fermion interpretation must resolve taste multiplicity, spin structure, and the particle/antiparticle assignment. Identifying $Q_{\rm mis}$ itself as a Dirac spinor would change its stated role.

\subsection{Proper time from orbital energy}

At $m>0$,
\begin{equation}
\partial_m\Omega_g=m(D_g^2+m^2I)^{-1/2}.
\end{equation}
The rate is one on harmonic rest modes and lies in $(0,1]$ at finite cutoff. With supplied positive lapse choose
\begin{equation}
H_{\rm matter}=N^{1/2}\Omega_gN^{1/2}.
\end{equation}
Its mass derivative remains positive even if $N$ and $D$ do not commute.

For the desired smooth metric symbol, set $E=\sqrt{h^{ij}p_ip_j+m^2}$. Then
\begin{equation}
H=NE,\qquad
\dot x^i=\frac{Nh^{ij}p_j}{E},
\qquad
\sqrt{N^2-h_{ij}\dot x^i\dot x^j}=\frac{Nm}{E}.
\end{equation}
The Legendre transform is
\begin{equation}
p_i\dot x^i-NE=-\frac{Nm^2}{E}
=-m\frac{d\tau}{dt}.
\end{equation}
Thus the orbital symbol reproduces $S_m=-m\int d\tau$. It agrees with the native rest insertion when the same spatial operator has the required metric symbol.

The rate lower bound decreases as the momentum cutoff increases: highly relativistic states have $m/E\to0$. Proper time remains positive along each massive timelike trajectory without a common lower bound over arbitrarily high momenta.

\subsection{An enumerated two-geometry transfer}

Let $g_0$ be the five-regular-tetrahedron boundary of a four-simplex. Let $g_1$ replace one tetrahedron by four regular tetrahedra. Their cochain dimensions are $30$ and $44$. The assigned regular metrics make them different physical geometries.

At $\alpha=7$, the restricted boundary-pair families have $1920$, $1992$, $1992$, and $2064$ causal histories. Summing their angular weights and dividing by one common scalar gives
\begin{equation}
A=
\begin{pmatrix}
0.11695083113895063&0.34220318033747954\\
0.34220318033747954&1
\end{pmatrix}.
\end{equation}
Its eigenvalues are approximately $-0.0001362342282$ and $1.11708706537$. Use the two-step transfer
\begin{equation}
K_g=A^\dagger A,\qquad
\operatorname{spec}K_g=
\{1.85597649216\times10^{-8},\,1.24788351161\}.
\end{equation}
This is strictly positive in the finite example. A common scalar normalization only changes the energy origin.

Choose
\begin{equation}
\Delta t=2a\sqrt{7+3/8}
=13.9484661297\,\tstar,\qquad
H_g=-\frac1{\Delta t}\log K_g,
\end{equation}
which gives, in inverse-tick units,
\begin{equation}
H_g=
\begin{pmatrix}
1.14086743423&-0.395787437728\\
-0.395787437728&0.119545050300
\end{pmatrix}.
\label{eq:Hgeo}
\end{equation}
The time interval is a declared synchronization to twice a $4$--$1$ boundary-normal duration. A moving particle can accumulate a different proper time.

Because $K_g$ is ill-conditioned, its logarithm is obtained by diagonalizing $A$ and taking twice the logarithm of the absolute eigenvalues. This avoids loss of relative precision in the small eigenvalue after numerical squaring.

\subsection{Normalized transport between physical geometries}

Stellar cochain interpolation supplies a rectangular map between the two spaces. Its maximum singular value in their physical inner products is $1.4832396974$. Inserting this raw map into the off-diagonal block of \eqref{eq:Hgeo} lowers the geometry-only marked ground energy by $0.1319919448$ relative to the vacuum. The effect is a normalization mismatch, much larger than $\eps$.

Let $P_0,P_1$ project onto the harmonic subspaces. Separate the raw map into harmonic and excited parts and take the rectangular polar factor on each full-rank sector. The combined transport $R$ obeys
\begin{equation}
R^\dagger R=I,\qquad RP_0=P_1R.
\label{eq:physicalR}
\end{equation}
This retains harmonic rest states and gives an isometry. Excited operators on the two distinct geometries need not intertwine under $R$.

Use
\begin{equation}
\widehat H_g=
\begin{pmatrix}
(H_g)_{00}I&(H_g)_{01}R^\dagger\\
(H_g)_{10}R&(H_g)_{11}I
\end{pmatrix},
\qquad
H(m)=\widehat H_g+\Omega_{g_0}(m)\oplus\Omega_{g_1}(m)
\label{eq:jointH}
\end{equation}
at unit reference lapse. The vacuum sector is explicitly $H_g$.

\begin{proposition}[Exact joint rest gap]
For \eqref{eq:jointH} with \eqref{eq:physicalR}, the marked and vacuum ground energies satisfy
\begin{equation}
E_{\rm marked,0}(m)-E_{\rm vacuum,0}=m.
\end{equation}
\end{proposition}
\begin{proof}
On the image of $R$, $\widehat H_g$ is unitarily equivalent to copies of $H_g$. The unmatched complement has energy $(H_g)_{11}$, which is at least the smallest eigenvalue $E_{g,0}$. Hence $\widehat H_g\ge E_{g,0}I$. Also $\Omega_g(m)\ge mI$. The transported harmonic sectors share rest energy $m$ and support a copy of the geometric ground eigenvector, attaining the lower bound $E_{g,0}+m$.
\end{proof}

The native value follows from the algebra, since a subtraction of eigenvalues of order one cannot resolve a $10^{-23}$ gap in double precision. At the control mass $m=0.23$, the numerical gap error is $4.45\times10^{-16}$.

\subsection{The joint clock and finite checks}

For a mass-independent geometric coupling and positive lapse, the complete generator has $\partial_mH\succeq0$. Equation \eqref{eq:integratedclock} supplies its integrated clock. The chosen packet both spreads on cochains and changes its geometry component coherently.

The vacuum transfer error is $4.45\times10^{-16}$, the orbital mass-derivative finite-difference error $2.84\times10^{-10}$, the integrated clock derivative error $9.90\times10^{-11}$, the composition error $2.15\times10^{-16}$, and the unitarity error $1.76\times10^{-15}$. The source derivatives are checked at a numerically resolvable control mass.

A local rephasing changes the basis and conjugates the kernel consistently. This verifies covariance under the supplied basis transformation. Electromagnetic link dynamics and its Gauss law would require further action terms and constraints.

The finite construction supplies a common generator for geometry, orbital position, and the clock. Its vacuum comes from an enumerated coarse transfer. Its rest-gap identity and positive clock survive the chosen geometry coupling. These statements are independent of whether the many-boundary model has the desired relativistic dispersion and interacting gravitational phase.

\section{Exact microscopic realization of a supplied matter operator}
\label{sec:refinement}

\subsection{Cochain interpolation and retraction}

Consider coarse and fine cochain complexes describing the same supplied coarse geometry. Let
\begin{equation}
P:C_c\to C_f,\qquad R:C_f\to C_c
\end{equation}
satisfy, degree by degree,
\begin{equation}
RP=I,\qquad d_fP=Pd_c,\qquad Rd_f=d_cR.
\label{eq:chainmaps}
\end{equation}
Here $P$ interpolates coarse cochains and $R$ integrates refined cochains back over coarse simplices. They are distinct from the normalized transport between physically different geometries in \eqref{eq:physicalR}.

Subcomplexes and commuting projections are standard structural ingredients of compatible discrete differential complexes \cite{feec}. The following choice of microscopic metric uses those identities to preserve the supplied coarse operator exactly.

\begin{theorem}[Positive metric completion]
\label{thm:metric}
Given positive metrics $M_c$ and $M_a$, define $Q_f=I-PR$ and
\begin{equation}
\boxed{M_f=R^\dagger M_cR+Q_f^\dagger M_aQ_f.}
\label{eq:Mf}
\end{equation}
Then $M_f$ is positive definite, $P$ is isometric from $(C_c,M_c)$ into $(C_f,M_f)$, and
\begin{equation}
C_f=\Ran P\oplus^\perp\ker R
\end{equation}
is preserved by $d_f$ and its metric adjoint.
\end{theorem}
\begin{proof}
Every $z\in C_f$ decomposes uniquely as $z=Px+y$ with $x=Rz$ and $y=Q_fz\in\ker R$. Since $Q_fP=0$ and $Q_fy=y$,
\begin{equation}
z^\dagger M_fz=x^\dagger M_cx+y^\dagger M_ay.
\end{equation}
Positivity follows, and the same equation proves isometry and orthogonality. The two chain-map identities in \eqref{eq:chainmaps} show that $d_f$ preserves each summand. In an orthogonal block decomposition its adjoint preserves the same summands.
\end{proof}

Let $J$ denote $P$ in orthonormal coordinates. The restricted adjoint on the image of $P$ is the transported coarse adjoint, so
\begin{equation}
J^\dagger J=I,\qquad D_fJ=JD_c.
\end{equation}
For every bounded Borel function on the finite spectrum,
\begin{equation}
f(D_f)J=Jf(D_c).
\end{equation}
In particular,
\begin{align}
\Omega_f(m)J&=J\Omega_c(m),\\
(\partial_m\Omega_f)J&=J(\partial_m\Omega_c),\\
e^{-it\Omega_f}J&=Je^{-it\Omega_c}.
\label{eq:functional}
\end{align}
The construction therefore preserves the complete dispersion and clock of the supplied operator.

\subsection{Locality and the interpretation of refined cells}

For a fixed finite subdivision block, the new metric couples cochains inside that block and through the already supplied coarse boundary interactions. It generally contains off-diagonal entries. Assigning the same cells an unchanged raw regular-simplex Hodge metric would define a different operator.

The theorem realizes a coarse orbital theory in a microscopic code subspace. It does not supply the coarse physical symbol or remove unwanted coarse species. A wrong speed or extra low-energy branch is carried into the code along with every correct feature.

Repeated refinements can preserve the chosen operator at each step. Direct and sequential interpolation must satisfy the appropriate composition maps. A refinement identity changes resolution within one selected construction; it does not prove that independent physical local time deformations obey the gravitational constraint algebra.

\subsection{Explicit stellar maps}

For a stellar insertion at the barycenter of one tetrahedron, let $b_v$ be the coarse barycentric coordinate vector of fine vertex $v$. A coarse oriented $p$-simplex $\sigma=(i_0,\ldots,i_p)$ interpolates to a fine simplex $\tau=(v_0,\ldots,v_p)$ with
\begin{equation}
P_{\tau\sigma}=\det[b_{v_a}^{\,i_b}]_{a,b=0}^p.
\end{equation}
The new vertex has coordinate $1/4$ at each vertex of its parent tetrahedron; old vertices retain their coordinate basis vectors.

For unchanged simplices, $R$ is the identity incidence. On the subdivided parent tetrahedron, $R$ adds the oriented child tetrahedra with signs fixed by the interpolation orientation. Direct multiplication checks all three identities in \eqref{eq:chainmaps}. Construct \eqref{eq:Mf}, factor it by Cholesky, and form $J$ in orthonormal coordinates.

Two refinements of the five-tetrahedron sphere give dimensions $30\to44\to58$. The saved calculation has direct-versus-sequential interpolation difference zero, Dirac intertwining error $5.99\times10^{-15}$, clock error $3.85\times10^{-15}$, and real-time error about $2.1\times10^{-14}$. Both original harmonic modes remain. The tested complementary Dirac spectrum has no extra zero mode and gap about $1.5081$.

\subsection{The full joint generator}

Apply Theorem \ref{thm:metric} separately to each physical geometry in \eqref{eq:jointH}. Lift the physical-geometry transport through the two resulting isometries. The coarse one-particle space has dimension $74$, and the fine one dimension $102$, with the same nine-state internal spectator factor.

On the code,
\begin{equation}
H_fJ=JH_c,\qquad C_fJ=JC_c,\qquad
U_f(t)J=JU_c(t),
\end{equation}
where $C=\partial_mH$. Transported source operators obey the same intertwining.

\begin{center}
\begin{tabular}{lr}
\toprule Joint check & Reference error\\\midrule
$J^\dagger J=I$ & $2.23\times10^{-15}$\\
Complete generator & $1.83\times10^{-13}$\\
Proper-clock operator & $7.79\times10^{-15}$\\
Real-time evolution & $3.11\times10^{-14}$\\
Two external test-source insertions & $1.83\times10^{-13}$\\
Code-invariance commutator & $2.99\times10^{-14}$\\
\bottomrule
\end{tabular}
\end{center}

At control mass $0.23$, complementary energies lie between about $1.6451$ and $3.7631$. These states do not mix with the code in the displayed generator. Their energies can overlap energies of excited code states. A global spectral separation above every code state is unnecessary for the stated exact invariance.

The two source insertions in this finite check are specified test operators. The intertwining proves preservation of their response. It does not substitute them for two calibrated held masses in a force experiment.

\section{The matter light cone and the retained limitations}
\label{sec:cone}

\subsection{The full Whitney metric test}

For barycentric coordinates $\lambda_i$ on one tetrahedron, the Whitney form associated with an oriented $p$-simplex $I=(i_0,\ldots,i_p)$ is
\begin{equation}
w_I=p!\sum_{r=0}^p(-1)^r\lambda_{i_r}
d\lambda_{i_0}\wedge\cdots\widehat{d\lambda_{i_r}}\cdots
\wedge d\lambda_{i_p}.
\end{equation}
The local metric is
\begin{equation}
(M_p)_{IJ}=\int_\sigma w_I\cdot w_J\,dV.
\end{equation}
Its barycentric second moments are
\begin{equation}
\int_\sigma\lambda_i\lambda_j\,dV
=\frac{V_\sigma}{20}(1+\delta_{ij}).
\end{equation}
Using these exact moments retains every off-diagonal entry and avoids a quadrature ambiguity.

On flat embedded periodic BCC and irregular Delaunay controls, the low-momentum speeds converge to one. Keeping the combinatorics but assigning every tetrahedron an intrinsic regular metric leaves a split. For the BCC test,
\begin{equation}
v_0=\sqrt{5/6},\qquad v_T=\sqrt{4/5},\qquad v_3=\sqrt{3/4}.
\end{equation}
The discrepancy therefore survives the full metric matrix. It is a property of the supplied intrinsic operator in that test.

\subsection{A periodic inverse normalization}

Choose an isotropic coarse metric with the same total intrinsic block volume before measuring dispersion. Match the metric of harmonic representatives to its exterior-form metric, then extend the positive congruence to all cochains and Bloch momenta.

For the BCC block the common target coordinate speed is $0.8908987181$; multiplication by the block length conversion $1.1224620483$ gives physical speed one. The irregular block has target coordinate speed $0.3110324586$ and conversion $3.2150985288$. No measured branch speed enters those normalization matrices.

At successively smaller diagonal momenta, the BCC relative spread is approximately
\begin{equation}
1.52358\times10^{-5},\quad3.80915\times10^{-6},\quad
9.52499\times10^{-7}.
\end{equation}
The irregular block gives
\begin{equation}
3.73828\times10^{-4},\quad1.86870\times10^{-4},\quad
9.34294\times10^{-5}.
\end{equation}
The leading cones approach the prescribed target. The visible irregular splitting decreases approximately linearly over the tested wave numbers. Its physical suppression and compatibility with arbitrary changing boundaries remain to be established.

\subsection{What one common interacting model still requires}

The finite joint operator, the periodic correction, and the geometric state-sum equality are compatible pieces with stated domains. The same interacting many-boundary construction must carry the desired spatial symbol and species interpretation, satisfy the necessary source and metric constraints, and reach the extended gravitational phase. These are substantive physical tests at the supplied cutoff.

The existence of a microscopic realization does not require uniqueness of its interactions. Its predictions do require fixing those interactions before testing the phase and source response. The refinement theorem supplies an exact way to implement a good coarse operator once that operator and its couplings are specified.

\section{Comparison with the main manuscript}
\label{sec:comparison}

\subsection{The current companion manuscripts}
This construction is incorporated in the current Astra manuscript
\cite{current}, in C.8, H.12--H.15, J.11--J.13, and Q.11--Q.11b.
The accounts use the same selected microscopic realization and distinguish
its finite identities from its unmeasured interacting phase and response.
The empirical companion \cite{empirical} reports the current galaxy
calculations, approximate cosmological fit, and progenitor-tagged run.
Those effective results do not establish the CDT phase or $G_R$.

\subsection{What is retained}

The capacity alphabet, injective four-port count, orientation doubling, closure law, and entropy are unchanged. The same marked generator gives $\eps$, and the same electron anchor calibrates $\tstar$. Geometry and capacity remain separate operator factors of one substrate. The retained position and mismatch fiber remain outside the replaceable capacity register.

The ordinary source couples to one metric. The capacity variable is read after reduction, consistently with Appendix N. Static density therefore couples through the constrained metric sector. The microscopic cut count is conditioned on boundary geometry, consistently with Appendix Q.11; an independent geometric degeneracy is not multiplied into the same entropy.

The event count has a fixed volume-per-event dictionary in the clock construction. A Monte Carlo move is an algorithmic step, a CDT slice is a foliation label, and a retained particle reads proper duration. The formalism keeps these three times distinct.

\subsection{What the continuation selects}

The coherent bond state is a global organization of the original local density matrix. The absolute angular rule fixes the scalar part of the angular generator. The paired-cell volume assignment specifies how clock events calibrate geometric cells. The connected patches and polar normalizations supply geometric interfaces with the desired coarse weight.

On the matter side, cochain placement follows the manuscript's proposed Dirac--Kähler route. The positive square root, physical-geometry polar transport, reference synchronization, and microscopic metric completion are explicit interaction choices. The finite tests establish their compatibility and consequences.

These choices define a realization within the broader ontology. The capacity count admits multiple such choices. Once they are fixed, the state-sum and operator identities follow mathematically.

\subsection{The distinct geometric normalizations}

In the matched scalar branch, Appendix C.5 fixes $G=\Gstar$ through the selected background matching $S_\infty=S_0$. Appendix Q.11 then gives $\Ach=4\log2\,\Lstar^2$. Under that matched face reading, J.10 uses $a\simeq2.530\Lstar$.

The paired-cell continuation supplies $v_*/\Lstar^4=(\pi/24)m_4^4$, assigns half that event volume per cell per tick, and reads $a$ from the regular-tetrahedron volume. At $m_4=2.35$ it gives $a=2.56812\Lstar$, $\Ach=2.85582\Lstar^2$, and $G_0=1.03002\Gstar$.

This is a substantive alternative geometric calibration, also recorded in current Astra C.8 and J.12. Its near agreement with the manuscript's matched value is useful, and its premises can be checked separately. The provisional chain constant and paired volume assignment prevent treating the displayed $1.03002$ as a high-precision parameter-free measurement of $G_R$.

The new regulator also changes the mixed-to-spatial channel ratio. The earlier $\alpha=11.312$ branch requires about $3.84049$ mixed Euclidean channels under literal face counting. The present $\alpha=7$ branch uses exactly three. The old runs continue to test their original branch.

\subsection{Appendix-by-appendix relation}

\begin{longtable}{@{}L{0.12\textwidth}L{0.39\textwidth}L{0.43\textwidth}@{}}
\toprule Main part & Relevant content & Relation to this construction\\\midrule\endhead
\midrule\multicolumn{3}{r}{\small Continued on next page}\\\endfoot
\bottomrule\endlastfoot
Abstract and introduction & Finite capacity spine, electron anchor, selected Einstein branch, open host phase. & The connected tensor maps and joint finite operator make selected interfaces explicit; the phase distinction is preserved.\\[4pt]
A & Units and canonical conventions. & $\Lstar=c\tstar$ and $\Gstar=c^3\Lstar^2/\hbar$ are retained. Event volume, host edge, and reference interval are separately specified.\\[4pt]
B & Positive face matching, $7+9$ split, injectivity, closure ensemble. & The exact native count and state are preserved. The bond encoding reorganizes existing routing-orientation states.\\[4pt]
C & Edge stiffness, diamond Green function, source ratio, physical $G=\Gstar$ matching. & The matching is distinguished from the angular/paired-volume bare $G_0$ and from a measured $G_R$.\\[4pt]
D & Weak-field bridge, electron scale, faithful renewal. & The native $\eps$ and scale survive. The ordinary Newton source is varied through the same metric parent.\\[4pt]
E & Transport and cosmological implementation. & No new transport or cosmological fit is inferred from the finite CDT correspondence.\\[4pt]
F & Spherical capacity reduction, boundary thermodynamics. & The static lapse observable is compatible with the one-metric source. Strong-field saturation dynamics remain separate.\\[4pt]
G & Many-Pasts, recorded histories, branch realization. & The retained/history split is respected; exported information is kept in the dilation.\\[4pt]
H.1--H.9 & Microscopic renewal and marked charged generator. & Capacity constants and rest energy are inherited with their stated generator prescription.\\[4pt]
H.9a & Relational geometry, maximal fusion, metric-Regge branch, proposed orbital placement. & Connected primitive interfaces realize a selected CDT weight. Cochain propagation is implemented and its speed limitations are measured.\\[4pt]
H.10 & Perfect-action embeddings, Einstein TT witness, geometric underdetermination. & The normalization and cochain theorem are concrete inverse realizations. Selection of the many-boundary phase remains physical.\\[4pt]
H.11 & Retained position, nine-state fiber, no-which-path criterion. & Branch-dependent reset schedules are tested; a fixed-background channel preserves the retained sector.\\[4pt]
I & Mass shells and gauge extensions. & The free orbital example leaves the internal mismatch factor intact. It does not derive the full species or chiral gauge theory.\\[4pt]
J & Capacity-decorated simulations, phase diagnostics, bare coupling conventions. & The $\alpha=7$ branch is specified independently. Earlier campaigns and registrations are not reassigned to it.\\[4pt]
K & Numerical and consistency checks. & New finite tests have fresh copied-workspace outputs and explicit tolerances.\\[4pt]
L & Fork accounting and historical calibration. & The provisional $m_4$ and selected interactions remain visible; the inverse $m_4$ target is not reported as measured.\\[4pt]
M & Saturated cosmological carrier. & No additional conclusion is drawn for this separate carrier model.\\[4pt]
N & One metric, source conservation, static and covariant reductions. & The Newton and lapse equations follow from the same parent before constraints are eliminated.\\[4pt]
O & Fixed event measure and source-clock release. & The fixed-event premise is distinguished from host simplex counting and from a volume-fixing algorithm.\\[4pt]
P & Record preservation, entropy and energy ledgers, history capacity. & Coherent renewal respects the retained factor. Global record-resource and differing-event-count questions remain.\\[4pt]
Q & Geometry-capacity factorization, lock, cut-area map, acceptance conditions. & The representation roles are retained. Connected Euclidean maps satisfy a concrete interface condition; complete Lorentzian and interacting audits remain.\\[4pt]
R & Reproduction code. & The present bundle adds a portable wrapper, reference inputs, regenerated outputs, and a separate verification report.\\
\end{longtable}

\subsection{Fixed event volume and unimodular reasoning}

Appendix O.1 makes the continuum premise that the event count supplies a fixed volume form and that the determinant is not an independent local metric degree of freedom. Under this premise, variation gives
\begin{equation}
R_{\mu\nu}-\frac14Rg_{\mu\nu}
=8\pi G\left(T_{\mu\nu}-\frac14Tg_{\mu\nu}\right).
\end{equation}
With $\nabla_\mu T^{\mu\nu}=0$, the Bianchi identity gives
$R+8\pi GT=4\Lambda$, and the ordinary Einstein equation follows with an integration constant $\Lambda$.

Fixing the total four-volume in a CDT simulation is a global restriction. By itself it does not fix the volume form pointwise or establish that no local determinant mode propagates after coarse-graining. The event-to-geometry map and spectrum must carry that stronger premise. Similarly, a source-clock release model determines a cosmological history only after its release law and initial data are specified.

\subsection{The precise status of the GR statement}

The selected metric action has Einstein equations, their Newton limit, and de Sitter solutions. The normalized microscopic state sum equals a selected CDT measure. The finite orbital extension carries its vacuum transfer, native rest gap, and a positive clock. Theorem \ref{thm:inherit} states how a compatible coarse phase and physical response become properties of the microscopic realization.

The current construction establishes the microscopic realization and the conditional route to Einstein behavior. The completed connected-host campaign adds stationary de~Sitter-like profiles and registered volume-stiffness agreement (Section~\ref{sec:alpha7results}). Start independence and the matter-source coefficient remain unresolved. The light-cone and record-coherence interfaces must also belong to that interacting model.

\section{Deferred matter-response calculation}

The registered trace measurement and subsequent scalar development leave the matter response unresolved. Further production is deferred pending a finite-geometry calibration that recovers known interactions and a power assessment on the calibrated observable. A future test must fix the geometry measure, volume prescription, topology and relational source preparation, establish stationarity, and check the physical conversions before comparing its response with $G_0$ and $\Gstar$.

If that same model has the required phase and response, the exact microscopic identities carry the result into its prepared code. A discrepancy identifies a failure or renormalization in the specified candidate. A newly selected interaction can define a new candidate, with its own predictions and test; it does not change the earlier measurement.

At low curvature the desired effective action has the form
\begin{equation}
\Gamma[g,\psi]
=\frac1{16\pi G_R}\int\sqrt{-g}(R-2\Lambda_R)
+\Gamma_m[g,\psi]+\Gamma_\partial+\sum_i c_i\mathcal O_i.
\end{equation}
Its leading local tensor structure, higher-curvature terms, matter spectrum, and physical boundary data require different observables. The finite realization fixes a concrete starting theory for those questions.

\section{Replication from a separate workspace}
\label{sec:replication}

\subsection{Bundle layout}

The replication directory contains a reference workspace with byte-identical copies of the required research inputs. Relative directories match the original layout. The manifest records each input's original path, copied path, byte count, and SHA-256.

\begin{Verbatim}[fontsize=\small]
replication/
  INPUT_MANIFEST.json
  run_checks.py
  verify_results.py
  requirements.txt
  reference/workspace/
    Opus carrier completion (math connections).tex
    Matter sector/
      dirac kahler/dk.py
      real complex/complexes.py
      chain_clock_2026-10-13/
        working_metric_prototype/
        microstructure_inverse/
        unified_micro_model/
        causal_primitive_realization_2026-10-16/
        clock_matter_bridge_2026-10-16/
        orbital_transfer_2026-10-16/
        force_test/intrinsic_probe.py
  runs/
    verification_02/
      RUN_REPORT.json
      logs/
      workspace/
\end{Verbatim}

The reference workspace also includes source reports and stored numerical fixtures. The calculation wrapper creates a fresh run workspace and executes the scripts there. The original research files and the reference inputs are not output destinations.

\subsection{Environment and commands}

The fresh verification used Python 3.13.2 on 64-bit Windows, NumPy 2.4.6, SciPy 1.17.1, and SymPy 1.13.3. The runner sets each numerical library to one thread and uses Windows Idle priority. The mathematical scripts also support ordinary Python path handling; this delivery's executed platform is Windows.

From the standalone directory, create an environment with the recorded dependencies and run:
\begin{Verbatim}[fontsize=\small]
python -m pip install -r replication/requirements.txt
python -B replication/run_checks.py --name my_replication
python -B replication/verify_results.py --name my_replication
\end{Verbatim}
The run name must be new. Existing run directories are retained, and the wrapper refuses to overwrite them. Each subprocess writes a log. A failed subprocess stops the sequence and leaves the completed outputs for diagnosis.

A targeted invocation can use the shipped dependency fixtures:
\begin{Verbatim}[fontsize=\small]
python -B replication/run_checks.py --stage lift --name lift_check
\end{Verbatim}
This is a check of the chosen stage against supplied inputs. Running the complete sequence recomputes the upstream blocks and matrices before the final lift.

The paper PDF can be rebuilt separately with pdfLaTeX through \texttt{build.py}. Scientific calculations do not depend on a TeX installation.

\subsection{Calculation order and dependencies}

\begin{longtable}{@{}L{0.15\textwidth}L{0.35\textwidth}L{0.44\textwidth}@{}}
\toprule Stage & Script & Main result and dependency\\\midrule\endhead
\midrule\multicolumn{3}{r}{\small Continued on next page}\\\endfoot
\bottomrule\endlastfoot
pair & \nolinkurl{block_map.py} & Builds the 49-to-13 normalized pair map.\\
exact\_parent & \nolinkurl{exact_parent.py} & Rational annihilator calculation using the pair code.\\
parent & \nolinkurl{parent_hamiltonian.py} & Positive degree-four parent and kernel certificate.\\
bonds & \nolinkurl{capacity_bonds.py} & Native density matrix, equivariant Fock encoding, and coherent cut factors.\\
renewal & \nolinkurl{renewal_event.py} & Coherent capacity renewal isometry.\\
angular & \nolinkurl{corner_model.py} & Shared replica/curvature rule, bare scale, parent pullback.\\
regge & \nolinkurl{regge_response.py} & Nonlinear variation of the selected metric action.\\
metric & \nolinkurl{metric_fourier.py} & Gauge identities, low-momentum tensor structure, same-source Newton residue.\\
volume & \nolinkurl{volume_state.py} & Compatible pair-volume state for the selected calibration.\\
frames & \nolinkurl{graph_frame.py} & Coherent moving-frame renewal on supplied graph states.\\
connected & \nolinkurl{connected_patch.py} & Four-cell mixed interface, paired-sector rank and normalization.\\
cycle & \nolinkurl{cycle_patch.py} & Six-cell nineteen-dimensional interface and held-out contraction checks.\\
slabs & \nolinkurl{slab_realization.py} & 1920-history weight identity, scalar source and sewing.\\
clock & \nolinkurl{marked_clock.py} & Native-energy intrinsic curves, duration sewing, free-line response.\\
coherence & \nolinkurl{renewal_coherence.py} & Reset counterexample and fixed-background repair.\\
positive\_\allowbreak clock & \nolinkurl{positive_clock_transfer.py} & Fréchet-log clock theorem checks and geometric-code lift.\\
\nolinkurl{boundary_clock} & \nolinkurl{boundary_clock_diagnostic.py} & Duration cumulants and coherent history compression.\\
whitney & \nolinkurl{whitney_orbital.py} & Full local metric matrices and flat/regular dispersion controls.\\
dressing & \nolinkurl{harmonic_dressing.py} & Periodic inverse normalization and residual splitting.\\
transfer & \nolinkurl{causal_boundary_transfer.py} & Enumerated two-boundary-state transfer at $\alpha=7$.\\
joint & \nolinkurl{joint_orbital_clock.py} & Normalized physical transport, exact rest gap, joint clock and motion.\\
cochains & \nolinkurl{cochain_block_completion.py} & Two-step cochain metric completion.\\
lift & \nolinkurl{joint_microscopic_lift.py} & Complete joint 74-to-102 operator and source lift.\\
\end{longtable}

The names in this table are relative to their source directories in the reference tree. The wrapper contains the complete relative paths. The computational scripts are preserved byte for byte. The wrapper supplies the copied working environment and controls output paths.

\subsection{Fresh verification}

All twenty-three stages completed with exit code zero in the fresh copied workspace. The summed subprocess runtime was $354.44$ seconds, about six minutes on the shared machine, with one numerical-library thread. Runtime depends on host load. The exact per-stage times and package versions are in \texttt{RUN\_REPORT.json}.

The following are reference values against which the verifier checks physical quantities and numerical residuals. Harmless floating-point differences in orthogonal bases are judged through invariant products, spectra, and residual norms.

\begin{longtable}{L{0.46\textwidth}L{0.46\textwidth}}
\toprule Quantity & Expected value or acceptance scale\\\midrule\endhead
Native closure parameter & $0.0298668443935226$\\
Native entropy & $7.4198000235709$\\
Electron rest gap & $4.185458795583485\times10^{-23}$\\
Paired-cell bare ratio at $m_4=2.35$ & $1.0300209831851719$\\
Connected target dimensions & $7$ and $19$\\
Causal slab count & $1920$\\
Primitive/direct log-weight discrepancy & Below $10^{-10}$\\
Boundary scalar sewing discrepancy & Below $10^{-10}$\\
Native curve coverage & Every declared history\\
Joint physical cochain dimensions & $30$ and $44$\\
Joint reference interval & $13.948466129734529$ ticks\\
Vacuum transfer and control rest gap & Errors below $10^{-11}$\\
Clock finite-difference derivative & Error below $10^{-7}$\\
Joint coarse and fine dimensions & $74$ and $102$\\
Full generator and source intertwining & Errors below $10^{-10}$\\
Clock and real-time intertwining & Errors below $10^{-10}$\\
Reference input hashes & Exact equality to manifest\\
\bottomrule
\end{longtable}

The tolerance scales are deliberately looser than the observed roundoff residuals. They test the claimed identities across ordinary numerical-library variation without requiring identical singular-vector signs. The physical scale inputs are compared separately from numerical algebra tolerances.

\subsection{Independent implementation recipes}

\paragraph{Capacity and scale.}
Enumerate four-element subsets of $\{-3,\ldots,3\}$. Compute \eqref{eq:K2}, solve \eqref{eq:closure}, and form the entropy with the $48$-fold routing-orientation multiplicity. Evaluate \eqref{eq:epsilon} with stable logarithms. Supply the declared $m_4$, then apply \eqref{eq:eventvol}, \eqref{eq:areascale}, and \eqref{eq:G0}. Record the value of every supplied input.

\paragraph{Geometric interfaces.}
Construct orthonormal intertwiner bases from Clebsch--Gordan coefficients using \eqref{eq:intertwiner}. Contract the specified chain or cycle with invariant edge metrics. Project external bundles with \eqref{eq:fusion3}. Apply the retained pair-code maps. Inspect singular values before inverting the Gram matrix; a map with deficient row rank cannot be repaired by normalizing it. Normalize with \eqref{eq:polar} and test both $F_eF_e^\dagger=I$ and the desired parent support.

\paragraph{Causal weight.}
Enumerate staircase slabs and deduplicate legal neighbors as sets of ordered simplex vertex tuples. Derive metric Gram matrices directly from the edge-length table. Compute each hinge area and dihedral angle independently of the tensor contraction. Include the half-boundary angle. Compare logarithms, since multiplying many small raw amplitudes obscures their relative discrepancy.

\paragraph{Clock.}
Fix initial marks independently of mass. Propagate a timelike tangent through intrinsic simplex charts, preserving its norm and future orientation. Add proper durations over segments. Evaluate log transfers at the actual $\eps$, and compute source derivatives analytically. A larger mass can serve as a finite-difference control with its role recorded separately.

\paragraph{Orbital operator.}
Build incidence matrices from oriented simplices and check $d^2=0$. Compute full Whitney matrices using exact barycentric moments. Cholesky-transform to orthonormal coordinates before comparing Hermiticity or spectra. Find harmonic projectors spectrally with a cutoff well below the nonzero gap. Normalize physical-geometry transport on harmonic and complementary sectors separately.

\paragraph{Microscopic cochain completion.}
Build barycentric interpolation and signed retraction. Verify \eqref{eq:chainmaps} before defining the metric. Form \eqref{eq:Mf}, check positivity, construct $J$, and evaluate $D_fJ-JD_c$. Functional-calculus identities can then be tested with the same matrices. When lifting the joint geometry coupling, apply the code separately to each physical geometry.

\subsection{Common replication failures}

Omitting orientation entropy changes the rest gap exponentially. Counting both directions of each undirected edge changes the stiffness convention. Adding raw shape norms to the already normalized CDT measure changes its geometry weights. Mixing the old and new $\alpha$ values changes the candidate. A fixed-length edge path gives an identically insensitive source. Subtracting eigenvalues to resolve the native $10^{-23}$ gap loses the signal. A direct matrix logarithm of an ill-conditioned squared transfer can corrupt its smallest eigenvalue. Each failure has a corresponding exact or numerical check in the sequence above.

The archived verifier scripts retain their original absolute-path provenance checks. The portable wrapper does not run those historical verifiers; it uses the bundle manifest and \texttt{verify\_results.py}. This keeps provenance auditing separate from reproducing physical calculations.


\subsection{Completed simulation evidence packet}

The original finite-construction replication bundle remains unchanged.
The additional directory \nolinkurl{evidence/20261020_alpha7_checkpoint}
preserves the registered connected-host protocol and amendment, all six
run configurations, the final verdicts and their analysis scripts, and
the separate trace registration and verdict. It also preserves the
scalar and variance-diagnostic readouts as development records.
\nolinkurl{INPUT_MANIFEST.json} identifies each original and copied file
by path, size and SHA-256. All eighteen alpha-seven frozen inputs were
verified before the manuscript update. The small evidence packet is
distinct from a complete raw-data distribution; the original profiles,
checkpoints and saved source geometries remain at the locations recorded
in the campaign files. No scientific simulation was rerun for this update.

\appendix
\section{Regge conventions and the CDT coupling map}
\label{app:couplings}

\subsection{Geometry from squared edge lengths}

The Euclidean construction can be reproduced without an embedding chosen by the original code. For a four-simplex with vertices $0,\ldots,4$, set $s_{ij}$ equal to its squared edge lengths and form
\begin{equation}
G_{ij}=\frac{s_{0i}+s_{0j}-s_{ij}}2,\qquad i,j=1,\ldots,4.
\end{equation}
The matrix is the Gram matrix of four edge vectors. Positive definiteness is the Euclidean simplex condition, and
\begin{equation}
V_4=\frac{\sqrt{\det G}}{4!}.
\end{equation}
For a triangle $(i,j,k)$,
\begin{equation}
A_{ijk}=\frac14
\sqrt{2(s_{ij}s_{ik}+s_{ij}s_{jk}+s_{ik}s_{jk})
-s_{ij}^2-s_{ik}^2-s_{jk}^2}.
\end{equation}
These formulas determine areas and volumes directly from the intrinsic metric.

One dihedral-angle algorithm factors $G$ to obtain vertex coordinates. For the hinge $(i,j,k)$, project the displacement vectors to the other two vertices onto the two-dimensional orthogonal complement of the hinge plane. Normalize the projected vectors to $u$ and $v$. The interior dihedral angle is $\theta=\arccos(u\cdot v)$. Clipping the scalar product to $[-1,1]$ removes roundoff excursions at the inverse-cosine call. Degenerate projections must be rejected before this step.

At unit spatial edge, the $(4,1)$ coordinates can be chosen as
\begin{equation}
x_0,\ldots,x_3=\frac1{\sqrt8}
\{(1,1,1),(1,-1,-1),(-1,1,-1),(-1,-1,1)\},
\end{equation}
all with fourth coordinate zero, and
$x_4=(0,0,0,\sqrt{\alpha-3/8})$.
For $(3,2)$ use an equilateral triangle of circumradius $1/\sqrt3$ in the first two coordinates and the two vertices
\begin{equation}
(0,0,\tfrac12,\sqrt{\alpha-7/12}),\qquad
(0,0,-\tfrac12,\sqrt{\alpha-7/12}).
\end{equation}
Every same-slice edge has squared length one and every mixed edge has squared length $\alpha$. Determinants give the two volume formulas in Section \ref{sec:geometric}. The Lorentzian normal heights follow from the corresponding continued metric; their positive proper durations contain $\sqrt{\alpha+3/8}$ or $\sqrt{\alpha+7/12}$.

\subsection{Reduction of the closed action to simplex counts}

For a closed periodic causal triangulation with the usual spherical spatial slices, the counting identities used here are
\begin{equation}
N_2^{\rm SL}=N_{41},\qquad
N_2^{\rm mix}=2N_0+N_{41}+2N_{32}.
\label{eq:countidentities}
\end{equation}
Both time orientations are included in $N_{41}$ and $N_{32}$. Boundaries require their explicit boundary terms; one should not apply \eqref{eq:countidentities} unchanged to an open slab.

Let $\bar A=A/a^2$, $\bar V=V/a^4$, and define
\begin{equation}
k=\frac{a^2}{8\pi G_0},\qquad
\bar\lambda=\lambda a^4,\qquad
\Sigma_i=\sum_{h\subset i}\bar A_h\theta_{hi}.
\end{equation}
The Euclidean curvature sign is
\begin{equation}
S_E=-k\sum_h\bar A_h\delta_h
+\bar\lambda\sum_\sigma\bar V_\sigma.
\end{equation}
For an Einstein cosmological term in these conventions,
$\lambda=\Lambda/(8\pi G_0)$. Inserting the deficits and the counting identities gives
\begin{equation}
S_E=c_0N_0+c_{41}N_{41}+c_{32}N_{32},
\end{equation}
with
\begin{align}
c_0&=-4\pi k\bar A_{\rm mix},\\
c_{41}&=-2\pi k(\bar A_{\rm SL}+\bar A_{\rm mix})
+k\Sigma_{41}+\bar\lambda\bar V_{41},\\
c_{32}&=-4\pi k\bar A_{\rm mix}
+k\Sigma_{32}+\bar\lambda\bar V_{32}.
\end{align}
The conventional count action is
\begin{equation}
S_E=-(\kappa_0+6\Delta_{\rm CDT})N_0
+\kappa_4(N_{41}+N_{32})
+\Delta_{\rm CDT}(2N_{41}+N_{32}).
\end{equation}
Equating coefficients yields
\begin{align}
\Delta_{\rm CDT}&=c_{41}-c_{32},\\
\kappa_4&=2c_{32}-c_{41},\\
\kappa_0&=-c_0-6(c_{41}-c_{32}).
\label{eq:couplingmap}
\end{align}
These equations provide an independent implementation of the coupling map.

The angular prescription fixes
\begin{equation}
\frac{a^2}{G_0}=\frac{16\log2}{\sqrt3},\qquad
k=\frac{2\log2}{\pi\sqrt3}.
\end{equation}
Consequently the dimensionless curvature couplings at fixed $\alpha$ do not depend on the provisional $m_4$. That constant calibrates the physical lattice scale and $G_0/\Gstar$. At $\alpha=7$ and $\bar\lambda=0$, evaluation of \eqref{eq:couplingmap} gives
\begin{equation}
\kappa_0=2.202053115156816,\quad
\Delta_{\rm CDT}=0.32613832803380927,\quad
\kappa_4=-0.30477461996038535.
\end{equation}
The independent volume values are
\begin{equation}
\bar V_{41}=0.07583447801333872,\qquad
\bar V_{32}=0.09140587903533463.
\end{equation}

\subsection{Physical volume and numerical volume control}

At $\alpha=7$, fixed $N_4=N_{41}+N_{32}$ does not imply fixed physical volume, since
\begin{equation}
V_4=a^4(\bar V_{41}N_{41}+\bar V_{32}N_{32})
\end{equation}
depends on the mixture of simplex types. A term $\epsilon(N_4-\bar N_4)^2$ restrains the total count. A term involving $V_4-\bar V_4$ restrains physical volume. Either can define a useful regulated ensemble, but their relative history weights differ.

Varying a physical $\bar\lambda$ at fixed $a$, $G_0$, and $\alpha$ moves all the conventional couplings:
\begin{align}
\frac{d\Delta_{\rm CDT}}{d\bar\lambda}
&=\bar V_{41}-\bar V_{32},\\
\frac{d\kappa_4}{d\bar\lambda}
&=2\bar V_{32}-\bar V_{41},\\
\frac{d\kappa_0}{d\bar\lambda}
&=-6(\bar V_{41}-\bar V_{32}).
\end{align}
Tuning only $\kappa_4$ at fixed $\kappa_0,\Delta_{\rm CDT}$ adds a chemical potential for $N_4$. Within an exactly fixed-$N_4$ sector it is constant. This operation is not generally the same trajectory as varying the physical volume coefficient. Replication of a phase result requires recording which prescription was used.

For an actual simulation, write the complete target probability before implementing moves:
\begin{equation}
\pi(T)\propto C_T^{-1}
\exp[-S_E(T)-S_{\rm control}(T)]\,W_{\rm dyn}(T).
\end{equation}
Here $W_{\rm dyn}$ includes any dynamical matter determinant or other retained geometry-dependent factor. For a proposal distribution $q(T'\mid T)$, the Metropolis--Hastings acceptance probability is
\begin{equation}
p_{\rm acc}(T\to T')=
\min\left\{1,\frac{\pi(T')q(T\mid T')}{\pi(T)q(T'\mid T)}\right\}.
\end{equation}
Proposal multiplicities, automorphism conventions, and forbidden-move rules must be included consistently. A hard volume window can suppress legal transitions near its boundary. Acceptance, autocorrelation, stationary distributions, and independent starts diagnose the sampler, while a correct local action comparison alone does not.

\section{Sources, Gaussian sewing, and a finite-duration force benchmark}
\label{app:sources}

\subsection{Integrating an external quadratic probe}

Let a scalar probe on a fixed triangulation have positive quadratic kernel $M_T$ after boundary conditions or zero-mode removal:
\begin{equation}
S_{\phi,T}=\frac12\phi^TM_T\phi-J^T\phi.
\end{equation}
Its Gaussian integral is
\begin{equation}
Z_{\phi,T}(J)=(2\pi)^{n/2}(\det M_T)^{-1/2}
\exp\left(\frac12J^TM_T^{-1}J\right).
\label{eq:gaussianprobe}
\end{equation}
The action inserted before integration and the determinant after integration are two descriptions of the same probe. Retaining only the exponential response while discarding the geometry-dependent determinant defines a quenched or externally conditioned calculation, whose status should be specified.

Separate the variables into boundary and interior values, $\phi=(\phi_B,\phi_I)$. With
\begin{equation}
M=
\begin{pmatrix}
M_{BB}&M_{BI}\\
M_{IB}&M_{II}
\end{pmatrix},
\end{equation}
integration over $\phi_I$ gives the Schur complement
\begin{equation}
M_{\rm eff}=M_{BB}-M_{BI}M_{II}^{-1}M_{IB},
\qquad
J_{\rm eff}=J_B-M_{BI}M_{II}^{-1}J_I,
\end{equation}
the determinant factor $(\det M_{II})^{-1/2}$, and the constant
$\tfrac12J_I^TM_{II}^{-1}J_I$ in $\log Z$.

When two slabs share the field $\phi_B$, their kernels and source vectors add on that interface. The next integration is another Schur complement. Associativity follows from ordinary Gaussian integration and
\begin{equation}
\det M=(\det M_{II})(\det M_{\rm eff}).
\end{equation}
This identity is the analytic basis of the finite scalar sewing check. It tests the source bookkeeping as well as the pure geometry contraction.

\subsection{The connected mass derivative}

For held trajectories with preparation independent of their masses, write
\begin{equation}
Z(m_1,m_2)=\sum_Tp_0(T)
\exp[-m_1\ell_1(T)-m_2\ell_2(T)],\qquad
\sum_Tp_0(T)=1.
\end{equation}
Direct differentiation gives
\begin{align}
\partial_{m_i}\log Z&=-\langle\ell_i\rangle_{m_1,m_2},\\
\partial_{m_1}\partial_{m_2}\log Z
&=\langle\ell_1\ell_2\rangle_{m_1,m_2}
-\langle\ell_1\rangle_{m_1,m_2}
\langle\ell_2\rangle_{m_1,m_2}.
\end{align}
At small masses the connected ratio therefore has expansion
\begin{equation}
\log\mathcal R=m_1m_2\Cov_0(\ell_1,\ell_2)
+O(m_1^2m_2,m_1m_2^2).
\end{equation}
Terms depending on one source alone cancel. The preparation must be differentiated if it depends on mass. Varying a mass-dependent path optimizer without including that derivative changes the observable.

If all lengths are in $\Lstar$, masses in $\hbar/(c\Lstar)$, and $T,r$ in tick-length units, the Newtonian target is
\begin{equation}
\frac{\Cov_0(\ell_1,\ell_2)}{T}
\longrightarrow\frac{G_R/\Gstar}{r}.
\label{eq:dimensionlessforce}
\end{equation}
This makes the dimensional conversion explicit. It separates a dimensionless response estimator from the physical value obtained after the electron calibration.

\subsection{A continuum finite-duration target}

A finite-duration test can be compared with the Gaussian Einstein limit before extrapolating $T$. The Euclidean massless Green function in four dimensions is
\begin{equation}
D_4(\tau,\mathbf r)=\frac1{4\pi^2(\tau^2+r^2)}.
\end{equation}
The two parallel static segments of duration $T$ have integral
\begin{align}
I(T,r)
&=\int_0^Tdt\int_0^Tds\,D_4(t-s,\mathbf r)\\
&=\frac1{2\pi^2}
\left[\frac{T}{r}\arctan\frac{T}{r}
-\frac12\log\left(1+\frac{T^2}{r^2}\right)\right].
\label{eq:finitegreen}
\end{align}
To derive it, write the double integral as
$2\int_0^T(T-u)D_4(u,\mathbf r)\,du$ and integrate
$1/(u^2+r^2)$ and $u/(u^2+r^2)$ separately.

For the usual linearized Einstein propagator and conserved nonrelativistic sources, each source vertex carries $m/2$, and the $44,44$ tensor residue is $16\pi G_R$. The connected cross term is consequently
\begin{equation}
\log\mathcal R_{\rm Gaussian}=4\pi G_Rm_1m_2\,I(T,r)
\end{equation}
in natural units. Since
\begin{equation}
\lim_{T/r\to\infty}\frac{I(T,r)}T=\frac1{4\pi r},
\end{equation}
the interaction tends to $-G_Rm_1m_2/r$.
At finite $T/r$ the deviation is known explicitly; a short interval should not be fit directly to an infinite-duration potential.

This benchmark uses the constrained Einstein response with a consistent gravitational Gaussian prescription. It is not a probability measure for the unrestricted conformal mode, and it does not replace the interacting history sum. Endpoint forces or apparatus contributions needed to hold the source are part of the conserved preparation. Finite box Green functions and curvature corrections replace \eqref{eq:finitegreen} when the relevant scales are comparable with the box or background curvature radius.

\section{Stable numerical normalization and the native mass scale}
\label{app:numerics}

\subsection{Full Gram normalization}

Given a complex $d$-by-$n$ map $R$ with $d\le n$, compute a thin singular-value decomposition
\begin{equation}
R=U\Sigma V^\dagger.
\end{equation}
The normalized map and its encoding are
\begin{equation}
F=UV^\dagger,\qquad I=VU^\dagger.
\end{equation}
They equal $(RR^\dagger)^{-1/2}R$ and its adjoint when all $d$ singular values are nonzero. This implementation avoids explicitly squaring the condition number. A singular-value cutoff is a rank decision, not a normalization convention. The retained ranks seven and nineteen must be verified before taking the polar factor.

Changing basis in the microscopic or coarse space conjugates the corresponding projectors. Comparisons should use $F F^\dagger$, $I I^\dagger$, singular values, and invariant contraction errors. Individual singular vectors can change sign or rotate inside degenerate singular subspaces across libraries while representing the same map.

The scalar raw packet norm is insufficient to fix the boundary operator. For two different packets $u,v$, sewing depends on
$u^\dagger RR^\dagger v$.
Normalizing just $R^\dagger u$ for one selected $u$ corrects that one vector but leaves the remaining cross terms. The complete Gram inverse fixes the entire target interface.

\subsection{The logarithm and its derivative}

Let $K=U\operatorname{diag}(\lambda_i)U^\dagger$ with $\lambda_i>0$. In that basis the Fréchet derivative is
\begin{equation}
\bigl[\mathcal D\log_K(C)\bigr]_{ij}
=L_{ij}C_{ij},\qquad
L_{ij}=
\begin{cases}
\displaystyle\frac{\log\lambda_i-\log\lambda_j}{\lambda_i-\lambda_j},
&\lambda_i\ne\lambda_j,\\[6pt]
1/\lambda_i,&\lambda_i=\lambda_j.
\end{cases}
\end{equation}
Near equal eigenvalues, use a stable divided difference or the continuous limit. The resolvent formula in Theorem \ref{thm:clock} proves positivity for positive $C$, including the noncommuting case. Replacing the derivative by $K^{-1}C$ would in general lose Hermiticity and give the wrong clock.

For the squared two-geometry transfer, its smallest eigenvalue is about $1.86\times10^{-8}$. Computing $\log(A^\dagger A)$ by first squaring can amplify relative roundoff in this eigenvalue. Here $A$ is Hermitian, so diagonalizing $A$ directly and using
$\log K=U\operatorname{diag}(2\log|\lambda_i(A)|)U^\dagger$
is stable. A zero eigenvalue would require restricting to a common support before defining the logarithm.

\subsection{Tiny native gaps and controlled tests}

The native electron gap is about $4.19\times10^{-23}$ in tick units. Double precision cannot resolve this by subtracting two order-one energies. The joint rest-gap theorem gives the exact analytic difference. Numerical checks at a control mass $0.23$ test the matrices and the implementation of that theorem.

Likewise, a rest amplitude with duration of order ten is
$e^{-\eps\ell}=1+O(10^{-22})$. Store its logarithm $-\eps\ell$ rather than subtracting its rounded exponential from one. A connected native source ratio of order $10^{-48}$ should be evaluated through its mass derivatives or cumulants.

These control masses are not changes in the physical electron calibration. They check analytic identities at a numerically resolved point. Predictions for the native mass must follow from the exact identity or a controlled expansion whose remainder is explicitly bounded.

The bare-scale sensitivity follows from
\begin{equation}
\frac{G_0}{\Gstar}=C\,m_4^{8/3},\qquad
C=0.10551926870131369,
\end{equation}
so
\begin{equation}
\frac{dG_0}{G_0}=\frac83\frac{dm_4}{m_4}.
\end{equation}
An exact target $G_0=\Gstar$ would require $m_4=C^{-3/8}$, approximately $2.324$. This is an inverse consistency target. It becomes evidence only when the chain constant is determined independently under the declared event law.

\section{Logical scope of the combined realization}
\label{app:scope}

\subsection{Existence, physical selection, and effective gravity}

An inverse construction can prove existence. Here the microscopic maps, selected simplex tensors, and source transport give a concrete member of the permitted class whose prepared code reproduces a chosen coarse theory. The construction makes additional dynamical commitments that can be examined as interactions on the substrate.

The original finite capacity measure does not determine every such commitment. The absolute angular offset, global bond coherence, paired volume assignment, normalized geometric vertices, and orbital metric must be stated with the resulting theory. Once fixed, they define predictions. A second microscopic realization with the same infrared action would not invalidate the first.

The following implications separate the mathematical and physical steps:
\begin{align}
&\text{native capacity state and chosen interactions}
\nonumber\\[-2pt]
&\hspace{1cm}\Longrightarrow
\text{normalized microscopic maps and bare curvature coefficient},
\\
&\text{maps, selected vertices, and matched history measure}
\nonumber\\[-2pt]
&\hspace{1cm}\Longrightarrow
\text{exact selected CDT generating functional},
\\
&\text{same ensemble in an Einstein regime with calibrated sources}
\nonumber\\[-2pt]
&\hspace{1cm}\Longrightarrow
\text{the corresponding microscopic effective GR and }G_R.
\end{align}
The first two implications are established for the specified construction and finite tests. The final implication is the inheritance theorem. Establishing its phase and response hypotheses for the new branch is the remaining physical calculation.

\subsection{The time requirement}

Three separate time constructions appear. The causal foliation labels allowed coarse simplices. The chain clock calibrates the physical event rate and $\tstar$. The finite transfer Hamiltonian uses a declared reference interval $\Delta t$, while a massive excitation accumulates its own positive integrated proper clock.

The local mass derivative and clock-composition identities have been checked. Equality of observables under different choices of spacetime foliation is stronger. Commuting disjoint replacement channels proves an order identity for those channels; it does not prove the algebra of overlapping gravitational constraints. A large CDT construction can instead supply an effective covariant regime if its observables restore the necessary symmetry. The specified finite matrices do not by themselves decide that limit.

At a zero of the lapse, an inverse-clock coordinate description can become singular. This does not imply that a supplied self-adjoint Hamiltonian ceases to generate unitary evolution in its reference parameter. The physical clock domain must be stated wherever lapse inversion is used. Expansion turning points and clock turning points are also distinct conditions.

\subsection{The deliverable and the physical conclusion}

The executable bundle reproduces the native finite constants, a positive paired parent, coherent bonds, connected geometric interfaces, the causal weight identity, native proper-time insertions, a finite joint operator, and an exact microscopic lift. It also retains counterexamples and limitations that define the usable domains of these constructions.

The selected Einstein parent supplies the classical Newton law and de Sitter solutions. The exact history identity makes a compatible coarse gravitational regime a regime of the microscopic realization. The current finite calculations have not measured the new ensemble's de Sitter phase or its interacting $G_R$. A matched simulation can test those properties without changing the microscopic definitions established here.

\begingroup
\small
\setlength{\parskip}{0pt}
\begin{thebibliography}{20}
\setlength{\itemsep}{3pt}
\bibitem{current} J. Chinitz, \emph{Entropic Scalar EFT: From Entanglement Microstructure to Gravity and Cosmic Structure}, current companion \texttt{Astra Quantum gravity completion (merged).tex}, 20 October 2026 evidence revision. The connected geometry and finite matter--clock construction appear in C.8, H.12--H.15, J.11--J.13, and Q.11--Q.11b; J.14 records the completed normalization and matter-response tests.
\bibitem{empirical} \emph{Empirical Results: History-dependent gravity and cosmological release}, companion \texttt{Empirical\_Results.tex}, 16 October 2026. Effective galaxy and cosmology comparisons; these do not measure the interacting response of the new CDT branch.
\bibitem{main} J. Chinitz, \emph{Entropic Scalar EFT: From Entanglement Microstructure to Gravity and Cosmic Structure}, workspace version \texttt{Opus carrier completion (math connections).tex}, 6 October 2026. The comparison and input manifest identify the exact file hash.
\bibitem{connected} \emph{Connected primitive patches for a causal CDT state sum}, research report, 16 October 2026. Replication reference: \nolinkurl{causal_primitive_realization_2026-10-16/CONNECTED_CDT_CONSTRUCTION.md}.
\bibitem{clock} \emph{Native rest transfer, causal duration, and coherent renewal}, research report, 16 October 2026. Replication reference: \nolinkurl{clock_matter_bridge_2026-10-16/CLOCK_MATTER_BRIDGE.md}.
\bibitem{orbital} \emph{Joint geometry orbital and clock completion}, research report, 16 October 2026. Replication reference: \nolinkurl{orbital_transfer_2026-10-16/JOINT_ORBITAL_COMPLETION.md}.
\bibitem{cdt} J. Ambj{\o}rn, J. Jurkiewicz, and R. Loll, \emph{Dynamically Triangulating Lorentzian Quantum Gravity}, Nucl.\ Phys.\ B \textbf{610}, 347--382 (2001), \href{https://arxiv.org/abs/hep-th/0105267}{arXiv:hep-th/0105267}.
\bibitem{ds} J. Ambj{\o}rn, A. G{\"o}rlich, J. Jurkiewicz, and R. Loll, \emph{The Nonperturbative Quantum de Sitter Universe}, Phys.\ Rev.\ D \textbf{78}, 063544 (2008), \href{https://arxiv.org/abs/0807.4481}{arXiv:0807.4481}.
\bibitem{feec} D. N. Arnold, R. S. Falk, and R. Winther, \emph{Finite element exterior calculus: from Hodge theory to numerical stability}, Bull.\ Amer.\ Math.\ Soc.\ \textbf{47}, 281--354 (2010), \href{https://arxiv.org/abs/0906.4325}{arXiv:0906.4325}.
\end{thebibliography}
\endgroup
\end{document}
